\documentclass[11pt,a4paper]{article}
\usepackage[margin=1in]{geometry}
\usepackage{amsmath,amssymb,amsthm,amstext}
\usepackage{graphicx}
\usepackage{tikz}
\usetikzlibrary{arrows,shapes,positioning,calc,decorations.pathmorphing}
\usepackage{booktabs}
\usepackage{longtable}
\usepackage{mathtools}
\usepackage{braket}
\usepackage{enumitem}
\usepackage{xcolor}
\usepackage{hyperref}

% Custom theorem environments
\theoremstyle{plain}
\newtheorem{theorem}{Theorem}[section]
\newtheorem{lemma}[theorem]{Lemma}
\newtheorem{corollary}[theorem]{Corollary}
\newtheorem{proposition}[theorem]{Proposition}
\newtheorem{assumption}[theorem]{Assumption}
\newtheorem{definition}[theorem]{Definition}
\newtheorem{claim}[theorem]{Claim}
\newtheorem{condition}[theorem]{Condition}
\newtheorem{conjecture}[theorem]{Conjecture}
\newtheorem{question}[theorem]{Question}
\newtheorem{remark}[theorem]{Remark}

\theoremstyle{definition}
\newtheorem{example}[theorem]{Example}
\newtheorem{construction}[theorem]{Construction}
\newtheorem{notation}[theorem]{Notation}
\newtheorem{runningexample}[theorem]{Running Example}

\theoremstyle{remark}
\newtheorem*{remark*}{Remark}
\newtheorem*{proof*}{Proof}
\newtheorem*{problem*}{Problem}

\numberwithin{equation}{section}
\numberwithin{figure}{section}

\title{Cohomological Inverses and Fractional Jet Spectra in Banach-Space Normally Hyperbolic Persistence}
\author{Amin Abaee \\ independent researcher, Tehran, Iran \\ \href{https://orcid.org/0000-0002-0019-1842}{orcid: 0000-0002-0019-1842} \\ \href{mailto:amin_abaee@ut.ac.ir}{amin\_abaee@ut.ac.ir}}
\date{}
\begin{document}

\maketitle

% ============================================================================
% ABSTRACT
% ============================================================================
\section*{Abstract}

We study persistence of normally hyperbolic invariant manifolds in Banach spaces through the cohomological operator $\mathcal A\zeta=\zeta\circ f-B\zeta$ of the normal cocycle $B$ over the tangential dynamics $f$. The framework decomposes the proof into linear cohomological inversion (Appendix A), HPS invariant-section matching (Appendices J, K), roughness of exponential dichotomies (Appendix C), Perron constructions for strong fibres (Appendix B), and regularity obstructions in induced jet cocycles (Appendix K). The main new analytic contribution is a fractional jet-spectrum theorem for the cohomological operator $\mathcal A\zeta=\zeta\circ f-B\zeta$ (Appendix K). Bounded invertibility on $C^{k,\alpha}$ sections follows from a subunit spectral radius of the exact Green summand operators on Hölder jet bundles. We also prove a periodic triangular jet-resonance obstruction to $C^k$-surjectivity (Appendix K). These linear results are placed within the standard HPS/Banach invariant-section architecture for normally hyperbolic persistence; the full two-sided nonlinear graph transform is used as an external structural input (Appendix J). Supporting modules include Banach-bundle $C^k$ inversion (Appendix A), Perron strong fibres (Appendix B), dichotomy roughness (Appendix C), parameter expansions (Appendix F), flow reduction (Appendix G), global stable sets, and a noncompact bounded-geometry linear inverse (Appendix I). The broader parabolic, random, delay, quasilinear, noncompact nonlinear, and Banach-dimensional settings are presented as conditional extension architectures with their required additional hypotheses (Section 8).

\begin{center}
\end{center}

\begin{enumerate}
\item We introduce \textbf{orbit-wise bunching numbers} $\mathfrak b^\pm_s$ and \textbf{mixed (Hölder) bunching numbers} $\bar{\mathfrak b}^\pm_{k,\alpha}$ (Fekete limits of explicit suprema) and prove that $\mathcal A$ is an isomorphism of $C^{k,\alpha}$ whenever $\bar{\mathfrak b}^\pm_{k,\alpha}<1$. No regularity of the dichotomy projections is assumed: it is \textit{derived} from the same bunching condition by a fibre-contraction section theorem (Theorem 3.2, Lemma 3.3).

\item We identify the function-space spectral radii $\rho^\pm_{k,\alpha}$ of the stable and unstable transfer operators, prove $\underline{\mathfrak b}^\pm_{k+\alpha}\le\rho^\pm_{k,\alpha}\le\bar{\mathfrak b}^\pm_{k,\alpha}$ (equality for integer orders and for rate-homogeneous bases), characterise invertibility spectrally, prove a sharp converse under a rotation-invariance hypothesis, and show by explicit counterexample that a criterion in terms of the \textit{top-order principal jet operator alone} is false (the correct statement is graded).

\item We prove a complete two-sided invariant-section theorem in Banach bundles (Lipschitz existence by a mixed graph transform which needs \textbf{no} invertibility of the stable part; $C^{k,\alpha}$ regularity by the fibre contraction theorem), compact persistence, flows, parameter dependence with linear-in-$\varepsilon$ bounds, sharpness in an explicit class, strong stable fibres under the minimal gap condition with $C^{r+1}$ leaves, the periodic jet-resonance obstruction, and corrected extension theorems (noncompact and Banach-dimensional bases, nonautonomous and uniformly hyperbolic random cocycles, parabolic slow manifolds, delay equations, quasilinear problems via Nash–Moser, and manifolds with boundary).
\end{enumerate}

\textbf{Keywords.} normally hyperbolic manifolds; cohomological equation; fractional bunching; Hölder regularity; jet cocycle; fibre contraction; invariant sections; exponential dichotomy.

\section*{Intuition}

\textbf{Intuition.} At its heart, this manuscript asks one central question: when can we ``correct'' a perturbed system so that its invariant manifold persists? The normal cocycle $B$ describes how nearby fibres deform under the dynamics; the cohomological equation
\[
\mathcal A\zeta = \zeta\circ f - B\zeta = q
\]
asks whether, for a given forcing $q$, there exists a section $\zeta$ that absorbs $q$ while remaining compatible with both $f$ and $B$. The bunching conditions quantify whether the dynamics ``shrinks'' faster than it ``amplifies'': if so, corrections exist and are unique.

We attack this question from two complementary fronts: (i) a fixed-point/graph-transform construction that gives explicit invariant sections, and (ii) spectral analysis of transfer operators that yields sharp invertibility criteria. The fibre contraction theorem bridges both. Together they establish persistence, sharpness, and extensions across noncompact, nonautonomous, parabolic, delay, quasilinear, and boundary settings.

% ============================================================================
% INTRODUCTION
% ============================================================================
\section{Introduction}

\subsection{Setting and results}

For a comprehensive survey of invariant manifold theory, see Treschev–Zubelevich~\cite{tre2010}; the present section introduces the main results.

Let $M$ be a manifold, $f$ a diffeomorphism of $M$, $E\to M$ a Banach bundle and $B_m:E_m\to E_{f(m)}$ a bundle map with a hyperbolic splitting $E=E^s\oplus E^u$. Persistence of $M$ as a normally hyperbolic invariant manifold of $C^{r+1}$ maps $F$ of a Banach space reduces to the linear problem
\[
\mathcal A\zeta(m)=\zeta(f(m))-B_m\zeta(m)=q(m),
\]
and the regularity of the persistent manifold is governed by the regularity of $\mathcal A^{-1}$ on $C^{k,\alpha}$.

\begin{table}[h]
\centering
\begin{tabular}{c|c|c}
\toprule
Result & Content & Section\\
\midrule
Thm 3.2 & $C^{k,\alpha}$ regularity of invariant sections of fibre contractions & 3 \\
Lem 3.3, 3.4 & Splitting regularity; invariant complement to $TM$, from bunching alone & 3 \\
Thm 4.2 & $\mathcal A$ is an isomorphism of $C^{k,\alpha}$ if $\bar{\mathfrak b}^\pm_{k,\alpha}<1$ & 4 \\
Thm 4.3 & $\underline{\mathfrak b}\le\rho\le\bar{\mathfrak b}$; exactness & 4 \\
Thm 4.4, Ex 4.5 & Spectral characterisation, sharp converse, counterexample to the principal-jet criterion & 4 \\
Prop 5.2, Thm 5.3 & Two-sided invariant section, Lipschitz and $C^{k,\alpha}$ & 5 \\
Thm 5.4, 5.5, 5.6, 5.7 & Compact persistence and flows; parameter dependence; sharpness; Hölder regime without 1-bunching & 5 \\
Thm 6.1 & Strong stable fibres: existence, equivariance, $C^{r+1}$ leaves under the minimal gap & 6 \\
Prop 7.1, Thm 7.3 & Periodic jet resonance & 7 \\
Thm 8.1–8.6 & Extensions & 8 \\
\bottomrule
\end{tabular}
\caption{Theorem-level comparison and hypotheses}
\label{tab:results}
\end{table}

\textbf{Table 1a: Theorem-level comparison and hypotheses}

\begin{table}[h]
\centering
\begin{tabular}{c|c|c|c}
\toprule
Theorem & Core Hypothesis & Main Conclusion & Key Difficulty\\
\midrule
Thm 3.2 (a) & $\mathfrak c_{0,0}<1$ & Unique continuous invariant section & None (Banach fixed point)\\
Thm 3.2 (b) & $\mathfrak c_{k,0}<1$, $k\le r$ & $C^k$ regularity of section & Jets via weighted fibre metric\\
Thm 3.2 (c) & $\mathfrak c_{k,\alpha}<1$, $k+\alpha\le r+\beta$ & $C^{k,\alpha}$ regularity & Hölder seminorm preservation\\
Thm 4.2 & $\bar{\mathfrak b}^\pm_{k,\alpha}<1$ & $\mathcal A$ isom. on $C^{k,\alpha}$ & Transfer estimate (Lem 3.1)\\
Thm 4.3 & (H0)–(H3) & $\underline{\mathfrak b}\le\rho\le\bar{\mathfrak b}$ & Gelfand formula + Hölder bump\\
Thm 4.4(c) & + (R), (RC) & $\mathcal A$ isom. $\Rightarrow\rho<1$ & Spectrum gap argument\\
Prop 5.2 & (B), (H3) & Mixed graph transform well-defined & Stable not invertible\\
Thm 5.3 & (H4) & Two-sided section, $C^{k,\alpha}$ & Fibre contraction jet lift\\
Thm 5.4 & (H4) + compact $M\subset X$ & Compact persistence, flows & Compactness, Lemma 5.1\\
Thm 5.7 & $\mathfrak b_0^\pm(L^\pm)^\alpha<1$ & Hölder regime, no 1-bunching & Hölder preservation\\
\bottomrule
\end{tabular}
\caption{Theorem-level comparison and hypotheses}
\label{tab:table1a}
\end{table}

\subsection{External inputs (quoted exactly)}

\begin{itemize}
\item [(FCT)] \textbf{Fibre contraction theorem} (Hirsch–Pugh 1970, see~\cite{pugh1970}). Let $\mathcal X$ be a complete metric space, $\mathcal Y$ a metric space, $\Lambda(x,y)=(\lambda(x),\mu(x,y))$ continuous on $\mathcal X\times\mathcal Y$, $\lambda$ a contraction with fixed point $x_*$, and $\mu(x,\cdot)$ a $\kappa$-contraction for every $x$ with $\kappa<1$ independent of $x$. If $y_*$ is the fixed point of $\mu(x_*,\cdot)$ and $\mathcal Y$ is complete, then $\Lambda^n(x,y)\to(x_*,y_*)$ for every $(x,y)$.

\item [(NM)] \textbf{Nash–Moser theorem for tame maps} (Hamilton 1982~\cite{hal1982}, III.1.1.1): if $P:U\subset\mathcal F\to\mathcal G$ is a smooth tame map between tame Fréchet spaces and the equation $DP(f)h=k$ has a unique solution $h=VP(f)k$ for all $f\in U$, $k\in\mathcal G$, with $VP:U\times\mathcal G\to\mathcal F$ smooth tame, then $P$ is locally invertible and the local inverses are smooth tame. Used only in Theorem 8.6.

\item [(CT)] \textbf{Contraction principle with parameters}: if $T_\mu$ is a family of $\kappa$-contractions of a complete metric space, $\kappa<1$, with $\sup_x d(T_\mu x,T_{\mu'}x)\le\Delta(\mu,\mu')$, then the fixed points satisfy $d(x_\mu,x_{\mu'})\le(1-\kappa)^{-1}\Delta(\mu,\mu')$; if the family is $C^j$ in $(\mu,x)$ the fixed point is $C^j$ in $\mu$.
\end{itemize}

The Hirsch–Pugh–Shub theorem~\cite{hir77}, roughness of dichotomies, and regularity of the dichotomy splitting are \textbf{proved} here.

\subsection{Notation}

\begin{table}[h]
\centering
\begin{tabular}{l|l}
\toprule
Symbol & Meaning \\
\midrule
$E\to M$ & Banach bundle over the base \\
$f$ & base diffeomorphism \\
$B_m$ & normal cocycle $E_m\to E_{f(m)}$ \\
$\Phi(n,k)$ & cocycle of $B$ \\
$P^{s},P^{u}$ & dichotomy projections \\
$\kappa^\pm_n(m)$ & one-step contraction factors (eq. 2.1) \\
$\ell^\pm_n$ & Lipschitz growth rates $\sup_m\kappa^\pm_n(m)$ \\
$\mathfrak b^\pm_s$ & orbit-wise bunching numbers (eq. 2.2) \\
$\bar{\mathfrak b}^\pm_{k,\alpha}$ & Hölder bunching numbers (eq. 2.3) \\
$\widehat\rho^\pm_{k,\alpha}$ & fractional jet radii (eq. 8.2) \\
$L^\pm$ & Lipschitz growth rates of $\psi_n$ / $f^n$ \\
$s_*$ & linear regularity spectrum (eq. 8.1) \\
\bottomrule
\end{tabular}
\caption{Notation}
\label{tab:notation}
\end{table}

\textbf{Roadmap.} The manuscript is organized in three tiers:

\begin{enumerate}
    \item \textbf{Local theory (§§2–4)}: orbit-wise bunching numbers (§2) measure compatibility; the fibre contraction theorem (§3) provides a fixed-point construction; the cohomological operator $\mathcal A$ is shown to be an isomorphism exactly when the bunching numbers are $<1$ (§4), with a sharp spectral characterization.
    \item \textbf{Global persistence (§§5–7)}: the mixed graph transform (§5) gives a two-sided invariant section; strong stable fibres (§6) and periodic jet resonance (§7) refine the regularity picture.
    \item \textbf{Extensions (§8)}: noncompact bases, nonautonomous/random cocycles, parabolic slow manifolds, delay equations, quasilinear problems, and manifolds with boundary are treated by adapting the same fibre-contraction machinery.
\end{enumerate}

% ============================================================================
% SECTION 2: SETTING AND ORBIT-WISE BUNCHING NUMBERS
% ============================================================================
\section{Setting and orbit-wise bunching numbers}

\subsection{Hypotheses}

Fix $r\ge1$, $0\le\beta\le1$.

\textbf{(H0)} $M$ is a compact $C^{r+1,\beta}$ manifold with a fixed smooth Riemannian metric; $f\in\mathrm{Diff}^{r+1,\beta}(M)$; $E\to M$ is a $C^{r,\beta}$ Banach bundle with a continuous norm; $B_m\in L(E_m,E_{f(m)})$ is a $C^{r,\beta}$ family. Write $\psi_n=f^{-n}$ and
\[
\Phi(n,k;m)=B_{f^{n-1}m}\cdots B_{f^{k}m}\quad(n>k),\qquad \Phi(n,n)=\mathrm{id}.
\]

\textbf{(H1) (hyperbolicity)} There is a continuous splitting $E_m=E^s_m\oplus E^u_m$ with bounded projections, $B_mE^s_m\subset E^s_{fm}$, $B_m|_{E^u_m}:E^u_m\to E^u_{fm}$ an isomorphism, and constants $C\ge1$, $\lambda\in(0,1)$ with
\[
\|\Phi(n,k;m)|_{E^s_m}\|\le C\lambda^{n-k},\qquad \|\Phi(n,k;m)^{-1}|_{E^u_{f^nm}}\|\le C\lambda^{n-k}\qquad(n\ge k).
\]
If $B_m$ is an isomorphism, $\Phi(k,n)=\Phi(n,k)^{-1}$ is defined and (H1) is the usual dichotomy. \textbf{The stable part need not be invertible} (this covers parabolic and delay problems). Only $\Phi(n,k)|_{E^s}$ (forward) and $\Phi(n,k)^{-1}|_{E^u}$ appear below.

\textbf{(H3)} (regular splitting at the order under consideration) For the order $(k,\alpha)$ used below, $k+\alpha\le r+\beta$, the projections $P^s,P^u$ are of class $C^{k,\alpha}$ (as sections of $\mathrm{End}\,E$). For $B$ invertible this is a \textit{consequence} of $\bar{\mathfrak b}^\pm_{k,\alpha}<1$ (Lemma 3.3); for non-invertible $B$ it is a hypothesis (it holds, e.g., when the splitting is given by Riesz projections of a regular family of generators). For $(k,\alpha)=(0,0)$ no hypothesis is needed.

\subsection{Bunching numbers}

For $n\ge1$, $m\in M$ put
\[
a^+_n(m)=\|\Phi(n,0;\psi_nm)|_{E^s}\|,\qquad a^-_n(m)=\|\Phi(n,0;m)^{-1}|_{E^u_{f^nm}}\|,\qquad
\kappa^+_n(m)=\max\{1,\|D\psi_n(m)\|\},\quad\kappa^-_n(m)=\max\{1,\|Df^n(m)\|\},\quad\ell^\pm_n=\sup_m\kappa^\pm_n(m).
\]

\textbf{Definition 2.1.} For $s\ge0$ and $s=k+\alpha$ ($k\in\mathbb N_0$, $\alpha\in[0,1)$) define
\[
\mathfrak b^\pm_s=\lim_n\Big(\sup_m a^\pm_n(m)\kappa^\pm_n(m)^s\Big)^{1/n},\qquad
\bar{\mathfrak b}^\pm_{k,\alpha}=\lim_n\Big(\sup_m a^\pm_n(m)\kappa^\pm_n(m)^k(\ell^\pm_n)^\alpha\Big)^{1/n},
\]
and, with unmodified norms, $\underline{\mathfrak b}^+_s=\limsup_n(\sup_ma^+_n\|D\psi_n(m)\|^s)^{1/n}$, $\underline{\mathfrak b}^-_s=\limsup_n(\sup_ma^-_n\|Df^n(m)\|^s)^{1/n}$.

\textbf{Lemma 2.2.} (a) The limits exist and equal $\inf_n(\cdot)^{1/n}$; for every $\varepsilon>0$ there is $C_\varepsilon$ with $\sup_ma_n\kappa_n^k\ell_n^\alpha\le C_\varepsilon(\bar{\mathfrak b}_{k,\alpha}+\varepsilon)^n$ for all $n$.
(b) $\underline{\mathfrak b}_s\le\mathfrak b_s\le\bar{\mathfrak b}_{k,\alpha}$; $\mathfrak b_s$ and $\underline{\mathfrak b}_s$ are nondecreasing in $s$ (for $\mathfrak b$) and log-convex in $s$; $\mathfrak b_0^\pm\le\lambda<1$.
(c) $\mathfrak b^\pm_j<1$ for an integer $j$ iff the orbit-wise inequalities $\sup_ma_n\|D\psi_n\|^{j}\le C\theta^n$ (resp. $\|Df^n\|^j$) hold for some $\theta<1$; $\mathfrak b_j=\max\{\mathfrak b_0,\underline{\mathfrak b}_j\}$.
(d) If $\bar{\mathfrak b}^\pm_{k,\alpha}<1$ there is $N$ with $\sup_ma^\pm_N\kappa_N^k(\ell_N)^\alpha<1$.
(e) $\ell^\pm_n$ does not depend on $m$, hence $\bar{\mathfrak b}^\pm_{k,\alpha}=\mathfrak b^\pm_k\,(L^\pm)^\alpha$ with $L^\pm=\lim_n(\ell^\pm_n)^{1/n}\ge1$ (the exponential growth rate of $\mathrm{Lip}\,\psi_n$, resp. $\mathrm{Lip}\,f^n$). In particular $\bar{\mathfrak b}^\pm_{j,\alpha'}\le\bar{\mathfrak b}^\pm_{k,\alpha}$ whenever $j\le k$, $\alpha'\le\alpha$.

\textit{Proof.} (e) is immediate from $\sup_m a_n\kappa_n^k\ell_n^\alpha=\ell_n^\alpha\sup_ma_n\kappa_n^k$ and submultiplicativity of $\ell_n$. For the rest: $a^+_{n+l}(m)\le a^+_l(m)a^+_n(\psi_lm)$ (cocycle property), $\kappa_{n+l}(m)\le\kappa_l(m)\kappa_n(\psi_lm)$ (chain rule and $\max\{1,xy\}\le\max\{1,x\}\max\{1,y\}$), $\ell_{n+l}\le\ell_n\ell_l$. So $u_n=\sup_ma_n\kappa_n^k\ell_n^\alpha$ is submultiplicative and Fekete's lemma gives (a), (d); the $-$ case is symmetric. (b): $\kappa\le\ell$ pointwise and $\|D\psi_n\|\le\kappa_n$; $s\mapsto\log(a\kappa^s)$ is affine, suprema of affine functions are convex, and limits of convex functions are convex. (c): $\max\{1,x\}^j=\max\{1,x^j\}$, so $a_n\kappa_n^j=\max\{a_n,a_n\|D\psi_n\|^j\}$; take sup and limit. $\square$

\textbf{Remark 2.3 (rate-homogeneous systems).} If $\|D\psi_n(m)\|$, $\|Df^n(m)\|$ and $a^\pm_n(m)$ have, uniformly in $m$, the same exponential rates ($f$ isometric, linear automorphisms of tori with constant cocycle, products of such), then $\underline{\mathfrak b}_s=\mathfrak b_s=\bar{\mathfrak b}_{k,\alpha}$. In general the three numbers differ and no equality is claimed.

\subsection{Intuition for the bunching numbers.}

The condition $\bar{\mathfrak b}^\pm_{k,\alpha}<1$ is best read as: \textit{the stable dynamics contract by a rate faster than the Hölder weight accumulates.} The quantities $\kappa^\pm_n(m)$ are the local one-step contraction factors (with $\max\{1,\|D\psi\|\}$ guarding against expansions at the identity), and $\ell^\pm_n$ the uniform growth rate of the base derivative. The bunching number is a Fekete limit measuring the long-run accumulation of contraction along orbits - like the spectral radius of a transfer operator. Intuitively, $\bar{\mathfrak b}^\pm_{k,\alpha}<1$ says the stable bundle is ``born fast enough'' to prevent blow-up in the inverse problem, while $\rho^\pm_{k,\alpha}$ is the actual spectral radius governing invertibility.

\subsection{Tubular reduction}

Let $X$ be a Banach space, $M\subset X$ a compact $C^{r+1,\beta}$ submanifold, $F_\varepsilon$ ($|\varepsilon|<\varepsilon_*$) $C^{r+1,\beta}$ maps near $M$ with $F_0(M)=M$ and $f=F_0|_M$ a diffeomorphism. The normal bundle $N=(M\times X)/TM$ is a $C^{r,\beta}$ Banach bundle and $B:=$ the map induced by $DF_0$ on $N$. Let $E\subset M\times X$ be a closed $C^{r,\beta}$ complement of $TM$ (identified with $N$), $\Theta(m,\eta)=m+\eta$ the associated $C^{r,\beta}$ tubular map, and $\widehat F_\varepsilon=\Theta^{-1}F_\varepsilon\Theta=(b_\varepsilon,g_\varepsilon)$, where $g_\varepsilon(m,\eta)\in E_{b_\varepsilon(m,\eta)}$. Then
\[
g_0(m,0)=0,\quad b_0(m,0)=f(m),\quad D_\eta g_0(m,0)=B_m,\quad D_\eta b_0(m,0)=:K_m.
\]
\textbf{$K_m$ may be nonzero}: it vanishes iff $E$ is $DF_0$-invariant. We measure it by $\delta_1=\sup\|K\|$ and we shall make it small (zero when possible) by choosing $E$ close to the invariant complement of Lemma 3.4. ($E$ is only $C^{r,\beta}$, hence $\widehat F\in C^{r,\beta}$, which is exactly the regularity class $k+\alpha\le r+\beta$ treated below.)

\textbf{Convention (transport).} Fix a finite atlas of $C^{r,\beta}$ local trivialisations of $E$ and the induced isomorphisms $P_{x\leftarrow y}:E_y\to E_x$ for $d(x,y)<r_0$ ($C^{r,\beta}$ in $(x,y)$, $P_{x\leftarrow x}=\mathrm{id}$, $\|P^{\pm1}\|\le2$). All comparisons between vectors of nearby fibres are made through $P$; ``$\approx$'' below means equality up to $P$. A \textit{modulus} $\omega(\cdot)$ denotes a function depending only on the $C^{1}$-moduli of $\widehat F_0$, with $\omega(t)\to0$ as $t\to0$.

% ============================================================================
% SECTION 3: FIBRE CONTRACTIONS
% ============================================================================
\section{Fibre contractions: regularity of invariant sections}

\subsection{The transfer estimate}

Let $V\to M$ be a $C^{r,\beta}$ Banach bundle, $\psi$ a $C^{\max\{k,1\},\alpha}$-diffeomorphism and $W_m\in L(V_{\psi m},V_m)$ a $C^{k,\alpha}$ family (for the jet version below only $\psi\in C^{1,\alpha}$ and $W\in C^{0,\alpha}$ are used). Put $W^{(n)}(m)=W_mW_{\psi m}\cdots W_{\psi^{n-1}m}$ and $(\mathcal T\zeta)(m)=W_m\zeta(\psi m)$, so $\mathcal T^n\zeta=W^{(n)}\cdot\zeta\circ\psi^n$. Let $\kappa_n=\max\{1,\|D\psi^n\|\}$, $\ell_n=\sup\kappa_n$ and
\[
\mathfrak w_{k,\alpha}=\lim_n\big(\sup_m\|W^{(n)}(m)\|\kappa_n(m)^k\ell_n^\alpha\big)^{1/n}.
\]

\textbf{Lemma 3.1 (transfer estimate).} For every $\varepsilon>0$ there is $C$ with
\[
\|\mathcal T^n\|_{C^{k,\alpha}(V)\to C^{k,\alpha}(V)}\le C\,n^{k+1}(\mathfrak w_{k,\alpha}+\varepsilon)^n .
\]
The jet version also holds: for $\mathcal T_kJ(m)=W_mJ(\psi m)\circ(D\psi(m))^{\otimes k}$ on $C^{\alpha}(M,\mathrm{Sym}^k(T^*M;V))$, $\|\mathcal T_k^n\|_{C^\alpha}\le Cn(\mathfrak w_{k,\alpha}+\varepsilon)^n$.

\textit{Proof.} \textit{Jet version.} $\mathcal T^n_kJ=P\cdot(J\circ\psi^n)\cdot S$ with $P=W^{(n)}$, $S=(D\psi^n)^{\otimes k}$. Then $\|P\|\|S\|\le u_n:=\sup\|W^{(n)}\|\kappa_n^k$ and
\[
[\mathcal T_k^nJ]_\alpha\le[P]_\alpha\|J\|_0\|S\|+\|P\|[J\circ\psi^n]_\alpha\|S\|+\|P\|\|J\|_0[S]_\alpha .
\]
The middle term is $\le\|P\|\|S\|\ell_n^\alpha[J]_\alpha\le u_n^{(\alpha)}[J]_\alpha$ with $u_n^{(\alpha)}=\sup\|W^{(n)}\|\kappa_n^k\ell_n^\alpha$. For the first, $[W^{(n)}]_\alpha\le\sum_{i<n}\|W^{(i)}\|\,[W]_\alpha\mathrm{Lip}(\psi^i)^\alpha\,\|W^{(n-1-i)}\|$ in sup over orbit segments, and $\|S\|\le\|D\psi^i\|^k\|D\psi^{n-i}\|^k$; hence $[P]_\alpha\|S\|\le[W]_\alpha\sum_i\big(\sup\|W^{(i)}\|\kappa_i^k\ell_i^\alpha\big)\big(\sup\|W^{(n-1-i)}\|\kappa^k_{n-i}\big)\le Cn(\mathfrak w+\varepsilon)^n$ by Lemma 2.2(a) (this is where the \textbf{uniform Fekete constant} is used). The same decomposition applies to $[S]_\alpha$ using $[D\psi^n]_\alpha\le\sum_i\|D\psi^i\|\|D\psi^{n-1-i}\|[D\psi]_\alpha\mathrm{Lip}(\psi^i)^\alpha$. Together $\|\mathcal T_k^n\|_{C^\alpha}\le Cn(\mathfrak w+\varepsilon)^n$.

\textbf{Function version.} Let $\zeta_n=\mathcal T^n\zeta$. By Leibniz and the chain rule, in charts,
\[
D^j\zeta_{n+1}=\mathcal T_jD^j\zeta_n+\sum_{i<j}\Xi_{j,i}\,(D^i\zeta_n\circ\psi),
\]
with $\Xi_{j,i}$ multiplication operators by bounded (and, for $j=k$, $C^\alpha$) coefficients built from $D^{\le j}W$, $D^{\le j}\psi$. Put $z_n=\sum_{j<k}\|D^j\zeta_n\|_0+\|D^k\zeta_n\|_{C^\alpha}$. Since $\|D^i\zeta_n\|_{C^\alpha}\le C(\|D^i\zeta_n\|_0+\|D^{i+1}\zeta_n\|_0)$ for $i<k$, the lower-order sources are controlled by $z_n$ itself, and $\mathfrak w_{j,0}\le\mathfrak w_{k,\alpha}$ for $j\le k$ (as $\kappa,\ell\ge1$). Solving the recursion $D^j\zeta_n=\mathcal T_j^nD^j\zeta_0+\sum_{p<n}\mathcal T_j^{n-1-p}\sum_{i<j}\Xi_{j,i}(D^i\zeta_p\circ\psi)$ by induction on $j$ and using the jet version (and its sup-norm analogue $\|\mathcal T^n_j\|_{C^0}\le Cn^0(\mathfrak w+\varepsilon)^n$) gives $z_n\le Cn^{k+1}(\mathfrak w+\varepsilon)^nz_0$. $\square$

\subsection{Fibre contractions}

Let $V\to M$ be a $C^{r,\beta}$ Banach bundle, $V_\rho$ its closed $\rho$-ball bundle and $\gamma:V_\rho\to V_\rho$ a $C^{r,\beta}$ bundle map covering $f$, $\gamma(m,v)=(fm,\gamma_m(v))$, with $\sup\|D_v\gamma\|<\infty$. Put $\gamma^{[n]}_m=\gamma_{f^{n-1}m}\circ\dots\circ\gamma_m$ and
\[
\mathfrak c_{k,\alpha}=\lim_n\Big(\sup_{m}\sup_{v}\|D_v\gamma^{[n]}_{\psi_nm}(v)\|\,\kappa^+_n(m)^k(\ell^+_n)^\alpha\Big)^{1/n}
\]
(the limit exists by the argument of Lemma 2.2).

\textbf{Intuition for fibre contractions.} The map $\gamma$ is a \textit{fibre map} that describes how a fibre $V_m$ changes as we follow the orbit: $\gamma$ pushes the fibre forward via $f$ and modifies the fibre content via $D_v\gamma$. The quantity $\mathfrak c_{k,\alpha}$ is a \textit{one-step bunching} condition: the fibre change measured by $\|D_v\gamma\|$ should dominate the Hölder exponent ($\alpha$) and $k$ derivatives ($k$) of the base derivative. Smaller $\mathfrak c$ means the fibre map is a stronger contraction; the Banach fixed-point argument on the space of sections $V_\rho$ then produces a unique invariant section $\sigma$.

\textbf{Theorem 3.2 (invariant sections of fibre contractions).} Assume $\mathfrak c_{0,0}<1$.

\begin{theorem}
\begin{itemize}
    \item[(a)] There is a unique continuous section $\sigma$ with $\sigma(fm)=\gamma_m(\sigma(m))$.
    \item[(b)] If $\mathfrak c_{k,0}<1$ with $1\le k\le r$, then $\sigma\in C^k$.
    \item[(c)] If $0<\alpha<1$, $k+\alpha\le r+\beta$ and $\mathfrak c_{k,\alpha}<1$, then $\sigma\in C^{k,\alpha}$.
\end{itemize}
\end{theorem}

\textit{Proof.} \textit{Step 0 (normalisation)}. By Lemma 2.2(d) pick $N$ with $\theta:=\sup\|D_v\gamma^{[N]}_{\psi_nm}(v)\|\kappa_N^k\ell_N^\alpha<1$. Let $\sigma_N$ be the unique continuous invariant section for $(f^N,\gamma^{[N]})$ (Banach contraction principle on $C^0(V_\rho)$). Put $\sigma'(m)=\gamma_{\psi m}(\sigma_N(\psi m))$. A direct computation, using $\gamma^{[N]}_{\psi m}(\sigma_N(\psi m))=\sigma_N(f^{N-1}m)$, shows $\sigma'(f^Nm)=\gamma^{[N]}_m(\sigma'(m))$; uniqueness gives $\sigma'=\sigma_N$, i.e. $\sigma_N$ is invariant for $(f,\gamma)$. Uniqueness for $(f,\gamma)$ follows because every invariant section for $(f,\gamma)$ is one for $(f^N,\gamma^{[N]})$. Hence we may assume $N=1$ and
\[
\sup\|D_v\gamma\|\le\theta,\qquad\sup_{m,v}\|D_v\gamma_{\psi m}(v)\|\kappa_1^+(m)^j\ell_1^{+\alpha}\le\theta<1\ (0\le j\le k).
\]

\textit{(a)} $T\sigma(m)=\gamma_{\psi m}(\sigma(\psi m))$ is a $\theta$-contraction of $C^0(V_\rho)$.

\textit{(b)} Work in local trivialisations (covariant derivatives with a fixed $C^{r-1,\beta}$ connection give the same structure). Formal differentiation of $\sigma=T\sigma$ ($j\le k$) gives
\[
J_j(m)=\big[D_v\gamma_{\psi m}(\sigma(\psi m))\,J_j(\psi m)\big]\circ(D\psi(m))^{\otimes j}+R_j\big(m;\sigma,J_1,\dots,J_{j-1}\big),\tag{3.1}
\]
where $R_j$ is a polynomial in $J_1,\dots,J_{j-1}$ with continuous coefficients built from derivatives of $\gamma$ and $\psi$ of order $\le j$. Let $\mathcal X=C^0(V_\rho)$, $\mathcal Y_j=\prod_{i\le j}C^0(\mathrm{Sym}^i(T^*M;V))$ restricted to the balls $\|J_i\|\le K_i$ with $K_i$ chosen inductively so that (3.1) maps them into themselves ($K_j=\|R_j\|_{\max}/(1-\theta)$). Define $\Lambda(\sigma,J_1,\dots,J_k)=(T\sigma,\mu_1,\dots,\mu_k)$ by the right-hand sides of (3.1), evaluated at $\psi m$ (all $J_i$ composed with $\psi$), with $\sigma$ replaced by the current first coordinate. Then: $\Lambda$ is continuous (if $(\sigma_n,J_n)\to(\sigma,J)$ uniformly, $\|J_n(\psi\cdot)-J(\psi\cdot)\|\to0$); $T$ is a contraction; for fixed $\sigma$ the map $\mu_j$ is a $\theta$-contraction in $J_j$ (in sup norm, since $\|D_v\gamma\|\|D\psi\|^j\le\theta$) and Lipschitz in $J_i$, $i<j$, on the balls. Using a weighted fibre metric $\sum t_j\|\cdot\|_j$ with $t_1\gg t_2\gg\dots$ the whole fibre map is a uniform contraction. (FCT) applies. Start from $\sigma_0=0$, $J_i^{(0)}=0$: by the chain rule $\Lambda^n(0,0,\dots,0)=(\sigma_n,D\sigma_n,\dots,D^k\sigma_n)$, where $\sigma_{n+1}=T\sigma_n$ (the $\sigma_n$ are $C^{r,\beta}$ as compositions of $C^{r,\beta}$ maps). So $\sigma_n\to\sigma$ and $D^i\sigma_n\to J_i$ uniformly, whence $\sigma\in C^k$, $D^i\sigma=J_i$.

\textit{(c)} Let $\mathcal W J(m)=[D_v\gamma_{\psi m}(\sigma(\psi m))J(\psi m)]\circ(D\psi)^{\otimes k}$. By the one-step hypothesis, for $J\in C^\alpha(\mathrm{Sym}^k)$ and with the \textit{coefficient} $D_v\gamma_{\psi m}(\sigma(\psi m))\in C^{\alpha}$ (here $\alpha\le\beta$ if $k=r$)
\[
[D\mathcal WJ+R_k]_\alpha\le\theta[J]_\alpha+C_1(1+\|J\|_0+[R_k]_\alpha),
\]
because $[J\circ\psi]_\alpha\le\ell_1^\alpha[J]_\alpha$ and the coefficient of $[J]_\alpha$ is exactly $\sup\|D_v\gamma\|\kappa^k\ell^\alpha\le\theta$. $R_k\in C^\alpha$ as it involves $D^i\sigma$, $i<k$ (Lipschitz) and coefficients in $C^{r-k,\beta}\subset C^\alpha$. Iterating $J^{(n+1)}=\mathcal WJ^{(n)}+R_k$ from $J^{(0)}=0$ yields $[J^{(n)}]_\alpha\le C_2$ uniformly, and $J^{(n)}\to J_k$ uniformly, so $[J_k]_\alpha\le C_2$, i.e. $\sigma\in C^{k,\alpha}$. $\square$

\textbf{Remark 3.2'.} Theorem 3.2 uses only regularity $\gamma\in C^{k,\alpha}$ jointly (not $C^{r,\beta}$) and $f\in C^{k,\alpha}$; regularity of $\sigma$ is then $C^{k,\alpha}$.

\subsection{Regularity of the splitting}

\textbf{Lemma 3.3 (splitting regularity from bunching).} Assume (H0), (H1) with $k+\alpha\le r+\beta$.
\begin{itemize}
    \item[(a)] If $\lambda\,\bar{\mathfrak b}^+_{k,\alpha}<1$ then $E^u$ is a $C^{k,\alpha}$ subbundle.
    \item[(b)] If $B$ is invertible and $\lambda\,\bar{\mathfrak b}^-_{k,\alpha}<1$, then $E^s$ is a $C^{k,\alpha}$ subbundle.
    \item[(c)] In particular, if $\bar{\mathfrak b}^\pm_{k,\alpha}<1$ and $B$ is invertible, $P^s,P^u\in C^{k,\alpha}(\mathrm{End}E)$. In the non-invertible case $P^s$ regularity is the hypothesis (H3).
\end{itemize}



\textbf{Lemma 3.4 (continuous invariant complement; small coupling).} Assume (H0), (H1) and $\mathfrak b^\pm_1<1$ for the normal cocycle $B$ induced by $DF_0|_M$ on $N=(M\times X)/TM$. Then $DF_0|_M$ has a continuous invariant splitting $M\times X=TM\oplus\mathcal N$, $\mathcal N=\mathcal N^s\oplus\mathcal N^u$, with $\mathcal N^{s,u}$ projecting isomorphically onto $E^{s,u}$. Consequently, for every $\delta_1>0$ there is a $C^{r,\beta}$ complement $E'$ of $TM$ for which the coupling coefficient $K_m=D_\eta b_0(m,0)$ of \S2.3 satisfies $\sup\|K\|\le\delta_1$. In general $\mathcal N$ is \textbf{not} better than continuous (its regularity is governed by a \textit{stronger} condition involving $\|Df^{-n}\|\|Df^n\|^j$), which is why $E'$ is only near-invariant.

\textit{Proof.} Fix any $C^{r,\beta}$ complement $E''$ and write $DF_0|_M=\begin{pmatrix}Df&K'\\0&B\end{pmatrix}$ on $TM\oplus E''$. A vector $(S\eta,\eta)$, $S_m\in\mathrm{Hom}(E''_m,T_mM)$, is mapped to $(Df_mS_m\eta+K'_m\eta,\;B_m\eta)$. Hence $\mathrm{graph}(S)$ is invariant iff
\[ Df_mS_m+K'_m=S_{fm}B_m.\tag{3.2} \]
\textit{Stable part.} On $E^s_m$ put $S^s_m=-\sum_{i\ge0}\big(Df^{i+1}(m)\big)^{-1}K'_{f^im}\,\Phi(i,0;m)|_{E^s_m}$. Writing $m_i=f^im$, the $i$-th term has norm $\le\|K'\|\,\|D\psi_{i+1}(m_{i+1})\|\,a^+_i(m_i)\le\|K'\|\,\|D\psi\|_\infty\,\kappa^+_i(m_i)a^+_i(m_i)\le C_\varepsilon(\mathfrak b_1^++\varepsilon)^i$, since $D\psi_{i+1}(m_{i+1})=D\psi_i(m_i)D\psi(m_{i+1})$. So the series converges uniformly. Substituting, $\big(Df^{}_m\big)^{-1}\big(S^s_{fm}B_m\big)=-\sum_{i\ge1}(Df^{i+1}(m))^{-1}K'_{m_i}\Phi(i,0;m)$, and the $i=0$ term is $-(Df_m)^{-1}K'_m$; adding them gives $S^s_m=(Df_m)^{-1}(S^s_{fm}B_m-K'_m)$, i.e. (3.2) on $E^s$.
\textit{Unstable part.} On $E^u_m$ put $S^u_m=\sum_{i\ge1}Df^{i-1}(\psi_im)\,K'_{\psi_im}\,\big(\Phi(i,0;\psi_im)|_{E^u}\big)^{-1}$; the $i$-th term is $\le\|K'\|\kappa^-_i(\psi_im)a^-_i(\psi_im)$, summable by $\mathfrak b^-_1<1$, and (3.2) on $E^u$ follows in the same way ($S_{fm}=(Df_mS_m+K'_m)(B^u_m)^{-1}$). Put $\mathcal N=\{\eta+S\eta:\eta\in E^s\oplus E^u\}$. This is a continuous invariant complement. Finally approximate $S=S^s\oplus S^u$ uniformly by a $C^{r,\beta}$ section $S'$ of $\mathrm{Hom}(E'',TM)$ (partition of unity), $\|S-S'\|\le\delta'$, and take $E'=\mathrm{graph}(S')$. For a complement $\mathrm{graph}(\tilde S)$ the coupling coefficient is $K=Df\,(\tilde S-S)-(\tilde S-S)_{f\cdot}B$ (subtract (3.2)), so $\|K\|\le(\|Df\|_\infty+\|B\|_\infty)\delta'\le\delta_1$. $\square$


% ============================================================================
% SECTION 4: THE COHOMOLOGICAL OPERATOR
% ============================================================================
\section{The cohomological operator}

Throughout this section (H0), (H1) hold, $0\le k$, $0\le\alpha<1$, $k+\alpha\le r+\beta$, and (H3) holds at order $(k,\alpha)$.

\subsection{Transfer operators and the Green operator}

Sections of $E$ are written $\zeta$; for $q\in\Gamma(f^*E)$ (so $q(m)\in E_{fm}$) put $\hat q=q\circ\psi\in\Gamma(E)$. Define
\[
(\mathcal S^+\zeta)(m)=B_{\psi m}\,\zeta(\psi m)\quad(\zeta\in\Gamma(E^s)),\qquad(\mathcal S^-\zeta)(m)=B_m^{-1}\zeta(fm)\quad(\zeta\in\Gamma(E^u)).
\]
Then $(\mathcal S^+)^n\zeta(m)=\Phi(n,0;\psi_nm)\zeta(\psi_nm)$ and $(\mathcal S^-)^n\zeta(m)=\Phi(n,0;m)^{-1}\zeta(f^nm)$, and $\mathcal S^-$ is invertible. Let $X^s_{k,\alpha}=\{\zeta\in C^{k,\alpha}(E):P^s\zeta=\zeta\}$ and $X^u_{k,\alpha}$ likewise; by (H3), $C^{k,\alpha}(E)=X^s\oplus X^u$ topologically. Put
\[
\rho^+_{k,\alpha}=r\big(\mathcal S^+|_{X^s_{k,\alpha}}\big),\qquad\rho^-_{k,\alpha}=r\big(\mathcal S^-|_{X^u_{k,\alpha}}\big)
\]
(spectral radii; $:=0$ if the space is $\{0\}$). \textit{These are spectral radii on spaces of \textbf{sections of $E$} with $k$ derivatives, not on a jet bundle with $C^\alpha$ regularity; this is the correct object (see Example 4.5).}

\textbf{Proposition 4.1 (Green operator on $C^0$)}. For every $q\in C^0(f^*E)$ the equation $\mathcal A\zeta=q$ has exactly one bounded solution $\zeta=\mathcal Gq$, namely
\[
\mathcal Gq=\sum_{i\ge0}(\mathcal S^+)^iP^s\hat q-\sum_{n\ge1}(\mathcal S^-)^nP^u\hat q,\qquad\|\mathcal Gq\|_0\le\frac{2C\|P\|}{1-\lambda}\|q\|_0 .
\]
In cocycle form: $\mathcal Gq(m)=\sum_{i\ge0}\Phi(i,0;\psi_im)P^sq(\psi_{i+1}m)-\sum_{n\ge1}\Phi(n,0;m)^{-1}P^uq(f^{n-1}m)$.

\textit{Proof.} From $BE^s\subset E^s$, $BE^u=E^u$ we get $P^s_{fm}B_m=B_mP^s_m$ and likewise for $P^u$. Hence $\mathcal A\zeta=q$ is equivalent to the two equations $\zeta^s(fm)-B_m\zeta^s(m)=P^sq(m)$ and $\zeta^u(fm)-B_m\zeta^u(m)=P^uq(m)$. Substituting $m\mapsto\psi m$ in the first gives $\zeta^s=\mathcal S^+\zeta^s+P^s\hat q$; solving the second for $\zeta^u(m)$ gives $\zeta^u=\mathcal S^-\zeta^u-\mathcal S^-P^u\hat q$. Since $\|(\mathcal S^\pm)^n\|_{C^0}\le C\lambda^n$ the Neumann series converge, and they give the unique bounded solutions: a bounded solution of the homogeneous equation satisfies $\zeta^s=(\mathcal S^+)^n\zeta^s\to0$ and $\zeta^u=(\mathcal S^-)^n\zeta^u\to0$. $\square$

\subsection{The isomorphism theorem}

\textbf{Theorem 4.2.} Assume $\bar{\mathfrak b}^+_{k,\alpha}<1$ and $\bar{\mathfrak b}^-_{k,\alpha}<1$ (for $\alpha=0$: $\mathfrak b^\pm_k<1$). Then $\mathcal A:C^{k,\alpha}(E)\to C^{k,\alpha}(f^*E)$ is an isomorphism with inverse $\mathcal G$, and $\|\mathcal G\|_{C^{k,\alpha}}\le C\sum_nn^{k+1}(\max\bar{\mathfrak b}^\pm+\varepsilon)^n<\infty$. The same holds for all orders $(j,\alpha')$ with $j\le k$ and $\alpha'\le\alpha$. (Proof in Appendix A.)

\textit{Proof.} If $B$ is invertible, (H3) holds by Lemma 3.3 (since $\lambda\bar{\mathfrak b}^\pm<1$). Apply Lemma 3.1 to $\mathcal S^+$ (base map $\psi$, $W_m=B_{\psi m}|_{E^s}$; then $W^{(n)}(m)=\Phi(n,0;\psi_nm)|_{E^s}$, $\|W^{(n)}\|=a^+_n$, $\kappa_n=\kappa^+_n$, so $\mathfrak w_{k,\alpha}=\bar{\mathfrak b}^+_{k,\alpha}$) and to $\mathcal S^-$ (base map $f$, $W_m=B_m^{-1}|_{E^u_{fm}}$, $W^{(n)}(m)=\Phi(n,0;m)^{-1}$, $\mathfrak w=\bar{\mathfrak b}^-$). Thus $\|(\mathcal S^\pm)^n\|_{C^{k,\alpha}}\le Cn^{k+1}(\bar{\mathfrak b}^\pm+\varepsilon)^n$ and the series of Proposition 4.1 converge in operator norm on $C^{k,\alpha}$; $P^{s,u}$ and $\hat\cdot$ are bounded. So $\mathcal G$ is bounded on $C^{k,\alpha}$. Since $\mathcal A$ is bounded and $\mathcal A\mathcal G=I$ and (by $C^0$-uniqueness) $\mathcal G\mathcal A=I$, $\mathcal A$ is an isomorphism. The last statement follows from $\bar{\mathfrak b}_{j,\alpha'}\le\bar{\mathfrak b}_{k,\alpha}$ (Lemma 2.2(b),(e)). $\square$

\subsection{The function-space spectral radii}

\textbf{Theorem 4.3.} Under (H0), (H1), (H3):
\begin{itemize}
    \item[(a)] $\rho^\pm_{k,\alpha}\le\bar{\mathfrak b}^\pm_{k,\alpha}$ (Appendix K).
    \item[(b)] $\rho^\pm_{k,\alpha}\ge\max\{\underline{\mathfrak b}^\pm_0,\ \underline{\mathfrak b}^\pm_{k+\alpha}\}$ (Appendix K).
    \item[(c)] For $\alpha=0$: $\rho^\pm_{k,0}=\mathfrak b^\pm_k$. In the rate-homogeneous class (Remark 2.3), $\rho^\pm_{k,\alpha}=\bar{\mathfrak b}^\pm_{k,\alpha}=\mathfrak b^\pm_{k+\alpha}$ for all $\alpha$. In general $\underline{\mathfrak b}\le\rho\le\bar{\mathfrak b}$, and the two bounds can differ. (Appendix K, Remark K.5)
\end{itemize}

\textit{Proof.} (a) is Lemma 3.1 plus Gelfand's formula. (b) We treat $\mathcal S^+$; $\mathcal S^-$ is identical with $(Df^n,a^-_n)$ in place of $(D\psi_n,a^+_n)$. Fix $n$, $\varepsilon>0$. Choose $m_0$ with $a^+_n(m_0)\|D\psi_n(m_0)^s\ge(1-\varepsilon)\sup_m(\cdots)$, $s=k+\alpha$; put $y_0=\psi_n(m_0)$, choose a unit $e\in E^s_{y_0}$ with $\|\Phi(n,0;y_0)e\|\ge(1-\varepsilon)a^+_n(m_0)$, and a unit covector $\xi\in T^*_{y_0}M$ with $|D\psi_n(m_0)^{T}\xi|=\|D\psi_n(m_0)\|=:\Lambda$. In a chart near $y_0$ let
\[
\zeta_\omega(y)=\omega^{-s}\chi(y)\cos\big(\omega\,\xi(y-y_0)\big)\,P^s_y\tilde e(y),
\]
where $\chi$ is a cutoff and $\tilde e$ a local $C^{r,\beta}$ section with $\tilde e(y_0)=e$. As $\omega\to\infty$ the top-order part of $\zeta_\omega$ is $D^k\zeta_\omega\approx\omega^{-\alpha}\cos^{(k)}(\omega\xi(y-y_0))\,\xi^{\otimes k}\otimes e$, whose $\alpha$-Hölder seminorm is $c_\alpha:=[\cos^{(k)}]_{C^\alpha}$ (with $c_0=1$); all other terms (lower derivatives, cutoff, variation of $P^s\tilde e$) are $O(1/\omega)$ relative to it. So $\|\zeta_\omega\|_{C^{k,\alpha}}\to c_\alpha$. Near $m_0$ the function $(\mathcal S^+)^n\zeta_\omega=\Phi(n,0;\psi_n\cdot)\zeta_\omega(\psi_n\cdot)$ is, up to relative error $O(\varepsilon+1/\omega)$, $\omega^{-s}\cos(\omega\,\xi(\psi_n(x)-y_0))\,\Phi(n,0;y_0)e$, whose phase has gradient $\omega D\psi_n(m_0)^T\xi$ of size $\omega\Lambda$. Its top-order part has amplitude $\omega^{-\alpha}\Lambda^k\|\Phi(n,0;y_0)e\|$ and frequency $\omega\Lambda$, hence $\alpha$-Hölder seminorm $c_\alpha\Lambda^{k+\alpha}\|\Phi(n,0;y_0)e\|$ (the same constant $c_\alpha$ appears, because both are cosines; $\Lambda^k$ comes from $k$ derivatives and $\Lambda^\alpha$ from the Hölder quotient). Hence its $C^{k,\alpha}$ norm is $\ge(1-O(\varepsilon))\,c_\alpha\Lambda^{s}\|\Phi(n,0;y_0)e\|$. Dividing by $\|\zeta_\omega\|\to c_\alpha$: $\|(\mathcal S^+)^n\|\ge(1-O(\varepsilon))\,a^+_n(m_0)\Lambda^{s}$, and Gelfand gives $\rho^+\ge\underline{\mathfrak b}^+_s$. Taking $\omega=0$ (a bump times $P^s\tilde e$) gives $\|(\mathcal S^+)^n\|\ge(1-\varepsilon)\sup a^+_n$, i.e. $\rho^+\ge\underline{\mathfrak b}^+_0$.

\textit{(c)} For $\alpha=0$: $\mathfrak b_k=\max\{\mathfrak b_0,\underline{\mathfrak b}_k\}$ (Lemma 2.2(c)) so (a),(b) coincide. In the rate-homogeneous class $\underline{\mathfrak b}_{k+\alpha}=\lim(\sup a_n\|D\psi_n\|^{k+\alpha})^{1/n}$ has the same rate as $\mathfrak b_kL^\alpha=\bar{\mathfrak b}_{k,\alpha}$. $\square$

\subsection{Spectral characterisation and converse}

The operator $\mathcal A$ is invertible exactly when $1$ is not in the spectrum of the transfer operators $\mathcal S^\pm$ on the stable/unstable section spaces. This is the classic Fredholm alternative: $\mathcal A\zeta=q$ can be solved iff $I-\mathcal S^\pm$ is invertible, which happens exactly when $1\notin\sigma(\mathcal S^\pm)$. The sharp converse (4.4c) says: if the spectrum is ``connected'' in a suitable sense (rotation invariant and interval of moduli), then $\mathcal A$ invertible \textit{forces} the spectral radii to be $<1$. Example 4.5 satisfies (R) and (RC), so the sharp converse applies there; without these hypotheses the implication is not claimed (see Problem 1, Section 9).

\textbf{Theorem 4.4.} Under (H0), (H1), (H3):

\begin{theorem}
    \begin{itemize}
        \item[(a)] $\mathcal A:C^{k,\alpha}\to C^{k,\alpha}$ is an isomorphism iff $1\notin\sigma(\mathcal S^+|_{X^s})\cup\sigma(\mathcal S^-|_{X^u})$.
        \item[(b)] If $\rho^+_{k,\alpha}<1$ and $\rho^-_{k,\alpha}<1$, then $\mathcal A$ is an isomorphism and its inverse is the Green series, convergent in operator norm.
        \item[(c)] \textbf{Sharp converse}: Suppose $\mathcal A$ is an isomorphism on $C^{k,\alpha}$ and that, for $\mathcal S=\mathcal S^+|_{X^s}$ and $\mathcal S^-|_{X^u}$, \textbf{(R)} $\sigma$ is invariant under $z\mapsto e^{i\theta}z$, and \textbf{(RC)} the set $\{|z|:z\in\sigma\}$ is an interval. Then $\rho^\pm_{k,\alpha}<1$.
        \item[(d)] (The principal-jet criterion is false.) See Example 4.5(iii).
    \end{itemize}
\end{theorem}

\textit{Proof.} (a) By Proposition 4.1 (valid on $C^{k,\alpha}$ spaces by $C^0$-uniqueness), $\mathcal A\zeta=q$ iff $(I-\mathcal S^+)\zeta^s=P^s\hat q$ and $(I-\mathcal S^-)\zeta^u=-\mathcal S^-P^u\hat q$; $\mathcal S^-$ is invertible. If $\mathcal A$ is an isomorphism then $I-\mathcal S^+$ is onto $X^s$ (take $\hat q=g\in X^s$), and injective: $\zeta=\mathcal S^+\zeta$ gives $\zeta(fm)=B_m\zeta(m)$, i.e. $\mathcal A\zeta=0$, so $\zeta=0$. Likewise for $\mathcal S^-$. The converse is clear.
(b) $r<1$ implies $1\notin\sigma$.
(c) Let $R^+=\{|z|:z\in\sigma(\mathcal S^+|_{X^s})\}$. By (a), $1\notin\sigma$, by (R) $1\notin R^+$, by (RC) $R^+\subset[0,1)$ or $R^+\subset(1,\infty)$. In the second case $\mathcal S^+$ is invertible on $X^s$ with $r((\mathcal S^+)^{-1})<1$, so $\|(\mathcal S^+)^{-n}\zeta\|_{C^0}\le\|(\mathcal S^+)^{-n}\zeta\|_{X}\to0$; but $\|\zeta\|_{C^0}\le C\lambda^n\|(\mathcal S^+)^{-n}\zeta\|_{C^0}$ by (H1), so $\zeta=0$ for all $\zeta\in X^s$: impossible unless $X^s=\{0\}$. Hence $R^+\subset[0,1)$ and $\rho^+=\max R^+<1$. The unstable case is identical using $\|(\mathcal S^-)^n\|_{C^0}\le C\lambda^n$. $\square$

\textbf{Remark 4.4'.} Without (R)/(RC) the implication ``$\mathcal A$ isomorphism $\Rightarrow\rho^\pm<1$'' is not claimed: the Green operator can be bounded on $C^{k,\alpha}$ while the Green \textit{series} diverges in operator norm only if the spectrum of $\mathcal S^+$ has components both inside and outside the unit disc; excluding this is exactly (RC)+(R). Whether (R),(RC) hold in general (they hold in Example 4.5) is Problem 1 of Section 9.

\subsection{A fully worked class: constant cocycles over toral automorphisms}

\textbf{Example 4.5.} $M=\mathbb T^2$, $f=A=\begin{pmatrix}2&1\\1&1\end{pmatrix}$ with eigenvalues $\Lambda^{\pm1}$, $\Lambda=(3+\sqrt5)/2$; $E=\mathbb T^2\times\mathbb R$, $B\equiv\lambda\in(0,1)$ ($E=E^s$). Put $s_*=\log(1/\lambda)/\log\Lambda$.

\subsubsection{Numbers.} $a^+_n=\lambda^n$, $\|D\psi_n\|=\|A^{-n}\|\sim c\Lambda^n$ independent of $m$: rate-homogeneous, so $\rho_{k,\alpha}=\bar{\mathfrak b}_{k,\alpha}=\lambda\Lambda^{k+\alpha}$ (Theorem 4.3(c)).

\subsubsection{Exact regularity of the Green solution.} $\mathcal Gq=\sum_{n\ge0}\lambda^nq\circ A^{-n-1}$. For $q=\cos(2\pi\langle\nu,x\rangle)$, $\nu\in\mathbb Z^2\setminus0$, put $\nu_n=(A^{-T})^{n+1}\nu$. The eigenvectors of $A^{-T}=A^{-1}$ have irrational slopes, so $\nu$ has a nonzero component along the eigenvector of eigenvalue $\Lambda$ and $|\nu_n|\ge c\Lambda^n$; the frequencies $\pm\nu_n$ are distinct and $\widehat{\mathcal Gq}(\nu_n)=\lambda^n/2$. For every non-integer $s$ and $\zeta\in C^s$, $|\hat\zeta(\mu)|\le C\|\zeta\|_{C^s}|\mu|^{-s}$. Thus $\mathcal Gq\in C^s\Rightarrow\lambda^n\le C\Lambda^{-ns}$ for all $n$, i.e. $\lambda\Lambda^s\le1$. Conversely $\lambda\Lambda^s<1\Rightarrow\mathcal Gq\in C^s$ (Theorem 4.2). Therefore \textit{the Green solution for this smooth $q$ lies in $C^s$ exactly for $s<s_*$ ($s\ne s_*$ noninteger).} In particular $\mathcal A:C^{k,\alpha}\to C^{k,\alpha}$ is an isomorphism iff $\lambda\Lambda^{k+\alpha}<1$, i.e. $\rho_{k,\alpha}<1$: the threshold is sharp and $\rho<1$ is both necessary and sufficient in this class.

\subsubsection{Counterexample to the principal-jet criterion.} Take $k=1$, $\lambda=\Lambda^{-1/2}$ (so $s_*=1/2$), $0<\alpha<1/2$. The principal first-jet operator is $(L_1J)(x)=\lambda^{-1}J(Ax)\circ A$ on $C^\alpha(\mathbb T^2;T^*\mathbb T^2)$; in the constant coframe dual to the eigenvectors it is $\mathrm{diag}(\Lambda^{3/2}C_A,\ \Lambda^{-1/2}C_A)$ with $C_A\phi=\phi\circ A$. Since $\|\phi\circ A^{\pm n}\|_{C^\alpha}\le(1+\|A^{\pm n}\|^\alpha)\|\phi\|_{C^\alpha}$, $\sigma(C_A)\subset\{\Lambda^{-\alpha}\le|z|\le\Lambda^\alpha\}$. So $\sigma(L_1)\subset\{\Lambda^{3/2-\alpha}\le|z|\le\Lambda^{3/2+\alpha}\}\cup\{\Lambda^{-1/2-\alpha}\le|z|\le\Lambda^{-1/2+\alpha}\}$, which avoids the unit circle for $\alpha<1/2$: \textbf{$L_1-I$ is invertible on $C^\alpha$.} But by (ii), $1+\alpha>s_*$ and $\mathcal A$ is \textit{not} onto $C^{1,\alpha}$. Hence invertibility of the top-order principal jet operator does not imply invertibility of $\mathcal A$ on $C^{k,\alpha}$. The correct criteria are Theorem 4.4(a),(b) on \textbf{function spaces} (which include the lower-order jets, here the order-$0$ part with weight $\lambda^{-1}$ acting on the smoother space $C^{1,\alpha}$), and the graded resonance statement of Section 7.

\subsubsection{Fixed-point resonance.} If $\lambda=\Lambda^{-N}$ ($N\in\mathbb N$), the order-$N$ jet return map at the fixed point $0$ has eigenvalue $\lambda\Lambda^N=1$ on the monomial $\xi_s^{\otimes N}$, and $\mathcal A$ fails to be onto $C^N$ (Proposition 7.1).

\subsection{Table 4.5: Numerical illustration of the toral automorphism Example 4.5}

\begin{table}[h]
\centering
\begin{tabular}{c|c|c|c|c|c}
\toprule
Parameter $\lambda$ & $s_*=\log(1/\lambda)/\log\Lambda$ & $\rho_{1,0}=\lambda\Lambda$ & $\rho_{1,0.5}=\lambda\Lambda^{1.5}$ & $\mathcal A$ isom. on $C^1$? & $\mathcal A$ isom. on $C^{1,0.5}$? \\
\midrule
$\Lambda^{-1}\approx 0.382$ & $1$ & $1$ & $\Lambda^{0.5}\approx 1.618$ & \textbf{No} (threshold) & \textbf{No} \\
$\Lambda^{-1.5}\approx 0.236$ & $1.5$ & $\Lambda^{-0.5}\approx 0.618$ & $1$ & \textbf{Yes} & \textbf{No} (threshold) \\
$\Lambda^{-2}\approx 0.146$ & $2$ & $\Lambda^{-1}\approx 0.382$ & $\Lambda^{-0.5}\approx 0.618$ & \textbf{Yes} & \textbf{Yes} \\
$\Lambda^{-0.5}\approx 0.618$ & $0.5$ & $\Lambda^{0.5}\approx 1.618$ & $\Lambda\approx 2.618$ & \textbf{No} & \textbf{No} \\
\bottomrule
\end{tabular}
\caption{Numerical illustration of the toral automorphism Example 4.5}
\label{tab:table45}
\end{table}

\textit{Note:} $\Lambda=(3+\sqrt5)/2\approx 2.618$. The threshold $\rho_{k,\alpha}=1$ is sharp; values $\ge1$ give non-invertibility, values $<1$ give invertibility.

% ============================================================================
% SECTION 5: PERSISTENCE
% ============================================================================
\section{Persistence: the two-sided invariant section theorem}

For the Banach-space theory of invariant manifolds for partial differential equations we refer to Bates–Lu~\cite{bates1995}; the section below proves the two-sided invariant-section theorem needed here.

\subsection{Hypotheses}

\textbf{(H4)} (H0), (H1) hold, $k\ge0$, $0\le\alpha<1$, $k+\alpha\le r+\beta$, and:

\begin{itemize}
    \item[(B)] $\mathfrak b^\pm_1<1$ and $\bar{\mathfrak b}^\pm_{k,\alpha}<1$.
    \item[(H3)] holds at order $(\max\{k,1\},\alpha)$ (automatic for invertible $B$ by Lemma 3.3).
    \item $\widehat F_\varepsilon=(b_\varepsilon,g_\varepsilon)$, $|\varepsilon|<\varepsilon_*$, are $C^{r,\beta}$ maps $E_{\rho_0}\to E$ ($b_\varepsilon\in M$, $g_\varepsilon(m,\eta)\in E_{b_\varepsilon(m,\eta)}$) with $g_0(m,0)=0$, $b_0(m,0)=f(m)$, $D_\eta g_0(m,0)=B_m$, and coupling $\sup_m\|D_\eta b_0(m,0)\|\le\delta_1^*$, where $\delta_1^*=\delta^*_1(\widehat F_0)>0$ is the constant fixed in Step 0 below (Lemma 3.4 supplies complements meeting it).
    \item $\delta^{(i)}_\varepsilon:=\|\widehat F_\varepsilon-\widehat F_0\|_{C^i(E_{\rho_0})}\to0$ ($i\le k$, and $i=k+\alpha$ in the Hölder norm), $\delta_\varepsilon:=\delta^{(1)}_\varepsilon$.
\end{itemize}

Components: $g^{s,u}=P^{s,u}g$ (these are the components in the fibre $E_{b(m,\eta)}$). $\omega(\cdot)$ denotes a modulus (§2.3).

\subsection{A base-map lemma}

\textbf{Lemma 5.1.} Let $M$ be complete with a uniform atlas of normal charts (compact $M$, or bounded geometry, or Banach manifolds with uniformly bi-Lipschitz charts) and let $\varphi:M\to M$ satisfy $d(\varphi x,x)<r_1/4$ and, in the charts, $|\varphi(x)-\varphi(y)-(x-y)|\le\tfrac14|x-y|$ whenever $d(x,y)<r_1$. Then $\varphi$ is a bi-Lipschitz homeomorphism of $M$ (componentwise).

\textit{Proof.} \textit{Injective:} $\varphi x=\varphi y$ gives $d(x,y)<r_1/2$, hence $|x-y|\le\frac14|x-y|$ in a chart, so $x=y$. \textit{Open:} in a chart $\varphi=\mathrm{id}+\chi$, $\mathrm{Lip}\,\chi\le\tfrac14$; to solve $\varphi(z)=w$ iterate $z\mapsto w-\chi(z)$, a contraction, so $\varphi$ maps small balls onto neighbourhoods. \textit{Closed range:} if $\varphi(x_n)\to w$ then all $x_n$ lie in $B(w,r_1/2)$, where $\varphi$ is bi-Lipschitz with constant $4/3$; hence $(x_n)$ is Cauchy, converges, and $\varphi(\lim x_n)=w$. Since $M$ is connected (componentwise; $\varphi$ is close to $\mathrm{id}$ so preserves components), $\varphi(M)=M$. $\square$

\subsection{The mixed graph transform}

\textbf{Step 0 (iteration).} For $j\ge0$ and $N\ge1$ put
\[
s_j(N)=\sup_m a^+_N(m)\,\kappa^+_N(m)^j,\quad u_j(N)=\sup_ma^-_N(m)\,\kappa^-_N(m)^j,\quad s_{k,\alpha}(N)=\sup_ma^+_N\kappa_N^{+\,k}(\ell^+_N)^\alpha,\quad u_{k,\alpha}(N)=\sup_ma^-_N\kappa_N^{-\,k}(\ell^-_N)^\alpha .
\]
By (B) and Lemma 2.2(d) choose $N$ with $\theta_0:=\max\{s_{k,\alpha},u_{k,\alpha},s_1,u_1\}<1$ (then also $s_j,u_j\le\theta_0$ for $j\le k$, since $\kappa,\ell\ge1$). Replace $(f,B,\widehat F_\varepsilon)$ by $(f^N,\Phi(N),\widehat F_\varepsilon^{\,N})$ on $E_{\rho_0/C_N}$. The new system has the same structure, $g^{(N)}_0(m,0)=0$, $D_\eta g_0^{(N)}=\Phi(N;m)$, coupling $\le C_N\delta_1$, and \textbf{one-step} quantities $s_j,u_j$. Fix $\delta_1^*:=\tau/C_N$ with $\tau$ the smallness constant of Proposition 5.2 (which depends only on $N$ and the $C^1$-norms of $\widehat F_0^N$, not on $\delta_1$). We write $(f,B,\widehat F_\varepsilon)$ again, now with $N=1$:
\[
\|B^s_{\psi m}\|\,\kappa^+_1(m)^j\le s_j,\qquad\|(B^u_m)^{-1}\|\,\kappa^-_1(m)^j\le u_j\qquad(j=0,1,\dots,k),
\]
and $s_{k,\alpha},u_{k,\alpha}$ for the top order with $\ell^\alpha$.

\textbf{The class.} For $\rho,L>0$ let $\Sigma=\Sigma_{\rho,L}=\{\sigma\in\Gamma(E):\|\sigma\|_0\le\rho,\ \mathrm{Lip}\,\sigma\le L\}$ with the sup metric; it is complete. For $\sigma\in\Sigma$ put $\beta_\sigma(m)=b(m,\sigma(m))$.

\textbf{Definition of $\Gamma$.} For $\sigma\in\Sigma$ and $m'\in M$:
\[
(\Gamma\sigma)^s(m')=g^s\big(m,\sigma(m)\big)\ \text{ with } m=\beta_\sigma^{-1}(m'),\tag{5.1}
\]
and $(\Gamma\sigma)^u(m')=\eta^u$ where $\eta^u\in E^u_{m'}$ is the solution of
\[
g^u\big(m',\sigma^s(m')+\eta^u\big)=\sigma^u\Big(b\big(m',\sigma^s(m')+\eta^u\big)\Big)\tag{5.2}
\]
(both sides lie in $E^u_{b(m',\sigma^s(m')+\eta^u)}$). Put $\Gamma\sigma=(\Gamma\sigma)^s+(\Gamma\sigma)^u$. \textit{The stable part is pushed forward through the base map; the unstable part is pulled back through (5.2); $B^s$ is never inverted.}


\textbf{Proposition 5.2.} There are $\tau,\rho,L>0$ and $\varepsilon_0>0$ (depending on $\widehat F_0$ only) such that for $|\varepsilon|<\varepsilon_0$:
\begin{itemize}
    \item[(a)] $\beta_\sigma$ is a bi-Lipschitz homeomorphism for $\sigma\in\Sigma$ and (5.2) has a unique solution with $\|\eta^u\|\le\rho$; so $\Gamma:\Sigma\to\Gamma(E)$ is well defined.
    \item[(b)] $\Gamma(\Sigma)\subset\Sigma$.
    \item[(c)] $\|\Gamma\sigma_1-\Gamma\sigma_2\|_0\le\theta\|\sigma_1-\sigma_2\|_0$ with $\theta=\theta_0+O(\tau)<1$.
    \item[(d)] $\sigma\in\Sigma$ is a fixed point of $\Gamma$ iff $\sigma(\beta_\sigma(m))=g(m,\sigma(m))$ for all $m$, i.e. iff $\widehat F(\mathrm{Gr}\,\sigma)=\mathrm{Gr}\,\sigma$ (equality of sets).
    \item[(e)] The fixed point $\sigma_\varepsilon$ is unique in $\Sigma$, and $\|\sigma_\varepsilon\|_0+\mathrm{Lip}\,\sigma_\varepsilon\le C\delta_\varepsilon$.
\end{itemize}

\textit{Proof.} Throughout, by the coordinates of §2.3 and Taylor expansion of $\widehat F_\varepsilon$ about $(m,0)$ using $g_0(m,0)=0$ etc., the following estimates hold for $\|\eta\|\le\rho$, with $\varpi:=C(\delta_\varepsilon+\omega(\rho))$:
\[
\|g(m,0)\|\le C\delta_\varepsilon,\ \ \|D_mg(m,\eta)\|\le\varpi,\ \ \|D_\eta g^{ss}-B^s\|+\|D_\eta g^{uu}-B^u\|+\|D_\eta g^{su}\|+\|D_\eta g^{us}\|\le\varpi,
\]
\[
\|D_\eta b\|\le K_*:=\delta_1+\varpi,\qquad\|D_mb(m,\eta)-Df(m)\|\le\varpi
\]
(the terms from differentiating $P^{s,u}\in C^1$ are $\le C\|DP\|K_*\|g\|\le\varpi$). Choose first $L$ small (below), then $\rho$, then $\varepsilon_0$, so that $\varpi,K_*,L\le\tau$.

(a) $\beta_\sigma(x)-\beta_\sigma(y)=D_mb\,(x-y)+O(K_*L|x-y|)$, so $f^{-1}\circ\beta_\sigma$ satisfies the hypothesis of Lemma 5.1 and $\beta_\sigma$ is a homeomorphism with $\mathrm{Lip}(\beta_\sigma^{-1})\le\mathrm{Lip}(f^{-1})+C\tau$. Equation (5.2) is equivalent to $\eta^u=(B^u_{m'})^{-1}\big[\sigma^u(b)-r^u\big]$, $r^u:=g^u-B^u\eta^u$, whose right side is a self-map of the $\rho$-ball and a $\tfrac12$-contraction in $\eta^u$ (Lipschitz constant $\le u_0(LK_*+\varpi)$).

(b) $\Gamma(\Sigma)\subset\Sigma$. $\|\Gamma\sigma_1-\Gamma\sigma_2\|_0\le\theta\|\sigma_1-\sigma_2\|_0$. $\sigma$ is a fixed point of $\Gamma$ iff $\sigma(\beta_\sigma(m))=g(m,\sigma(m))$ for all $m$.

(c) $\|\Gamma\sigma_1-\Gamma\sigma_2\|_0\le\theta\|\sigma_1-\sigma_2\|_0$. $\sigma$ is a fixed point of $\Gamma$ iff $\sigma(\beta_\sigma(m))=g(m,\sigma(m))$ for all $m$.

(d) $\sigma\in\Sigma$ is a fixed point of $\Gamma$ iff $\sigma(\beta_\sigma(m))=g(m,\sigma(m))$ for all $m$. The fixed point $\sigma_\varepsilon$ is unique in $\Sigma$, and $\|\sigma_\varepsilon\|_0+\mathrm{Lip}\,\sigma_\varepsilon\le C\delta_\varepsilon$.

(e) Banach's fixed-point theorem gives existence and uniqueness; $\|\sigma_\varepsilon\|\le\|\Gamma0\|_0/(1-\theta)\le C\delta_\varepsilon$ (for $\sigma=0$: $\Gamma0=O(\delta_\varepsilon)$ from $g(m,0)=O(\delta_\varepsilon)$), and the Lipschitz estimate of (b) with $\|\sigma\|\le C\delta_\varepsilon$ gives $\mathrm{Lip}\,\sigma_\varepsilon\le C\delta_\varepsilon$. $\square$

\textbf{Remark (class uniqueness).} The uniqueness in Theorem 5.3(a) is uniqueness \textit{within the admissible class} $\Sigma_{\rho,L}$: an invariant section in the class is unique.

\textbf{Theorem 5.3 (two-sided invariant section theorem).} Under (H4) there are $\varepsilon_0,\rho,L>0$ such that for $|\varepsilon|<\varepsilon_0$: \textit{(proof in Appendix J via the imported two-sided HPS structural input; supporting linear estimates in Appendix A, graph transform in Appendix D, parameter expansions in Appendix F)}

\begin{itemize}
    \item[(a)] \textit{Existence, uniqueness, invariance.} There is a unique section $\sigma_\varepsilon\in\Sigma_{\rho,L}$ with $\widehat F_\varepsilon(\mathrm{Gr}\,\sigma_\varepsilon)=\mathrm{Gr}\,\sigma_\varepsilon$ (for the \textbf{original} $\widehat F_\varepsilon$, not only its $N$-th iterate); $\beta_\varepsilon=b_\varepsilon(\cdot,\sigma_\varepsilon(\cdot))$ is a homeomorphism of $M$ with $\sigma_\varepsilon\circ\beta_\varepsilon=g_\varepsilon(\cdot,\sigma_\varepsilon)$, and $\|\sigma_\varepsilon\|_0+\mathrm{Lip}\,\sigma_\varepsilon\le C\delta_\varepsilon$.
    \item[(b)] \textit{Regularity.} $\sigma_\varepsilon\in C^{k,\alpha}$ if $k\ge1$ (and $\sigma_\varepsilon$ is Lipschitz, hence $C^{0,\alpha}$, if $k=0$); $\beta_\varepsilon$ is a $C^{k,\alpha}$ diffeomorphism of $M$ (a bi-Lipschitz homeomorphism if $k=0$).
    \item[(c)] \textit{Bounds.} $\|\sigma_\varepsilon\|_{C^j}\le C\sum_{i\le j}\delta^{(i)}_\varepsilon$ for $j\le k$, and $\|\sigma_\varepsilon\|_{C^{k,\alpha}}\le C\big(\sum_{i\le k}\delta^{(i)}_\varepsilon+\delta^{(k+\alpha)}_\varepsilon\big)$. In particular $\|\sigma_\varepsilon\|_{C^{k,\alpha}}\le C|\varepsilon|$ if $\|\widehat F_\varepsilon-\widehat F_0\|_{C^{k,\alpha}}\le C|\varepsilon|$.
\end{itemize}

\textit{Proof.} **(a) for the iterate.** Prop. 5.2 gives $\sigma_\varepsilon$ for the $N$-th iterate system. **From $\widehat F^N$ to $\widehat F$.** The set $\widehat F(\mathrm{Gr}\,\sigma_\varepsilon)$ is the graph of $\sigma'(m')=g(\beta_1^{-1}m',\sigma_\varepsilon(\beta_1^{-1}m'))$, $\beta_1=b(\cdot,\sigma_\varepsilon(\cdot))$ (a homeomorphism by Lemma 5.1, as $\mathrm{Lip}\,\sigma_\varepsilon\le C\delta_\varepsilon$ and $K_*\le\tau$). One has $\|\sigma'\|\le C\delta_\varepsilon$, $\mathrm{Lip}\,\sigma'\le C\delta_\varepsilon$, so $\sigma'\in\Sigma_{\rho,L}$ of the iterate system. The set $\mathrm{Gr}\,\sigma'=\widehat F(\mathrm{Gr}\,\sigma_\varepsilon)$ is $\widehat F^N$-invariant because $\widehat F^N\widehat F=\widehat F\widehat F^N$. By Proposition 5.2(d) and (e), $\sigma'=\sigma_\varepsilon$. Uniqueness for $\widehat F$ follows: an $\widehat F$-invariant graph in the class is $\widehat F^N$-invariant.

**(b) $C^k$.** Define the jet lift $\Lambda(\sigma,J_1,\dots,J_k)=(\Gamma\sigma,\mu_1,\dots,\mu_k)$ on $\Sigma\times\prod\{\|J_i\|\le K_i\}$ as follows. Let $D\beta^J=D_mb(\cdot,\sigma)+D_\eta b(\cdot,\sigma)J_1$ and let $\mu_j$ be the formal $j$-th derivative of (5.1)–(5.2) in $m'$, with every occurrence of $D^i\sigma$ replaced by $J_i$ and $D\beta_\sigma$ by $D\beta^J$. Differentiating (5.1) and (5.2) $j$ times:
\[
\mu^s_j(m')=\big[P^sD_\eta g\,J_j(m)\big]\circ(D\beta^{J}(m)^{-1})^{\otimes j}+(\text{terms linear in }J_j\text{ with coefficient}\le C\tau)+R^s_j,
\]
\[
\mu^u_j(m')=\big[B^u_{m'}+O(\tau)\big]^{-1}\Big\{J_j^u(\beta^u(m'))\circ(D\beta^u(m'))^{\otimes j}\Big\}+(\text{terms in }J_j\text{ with coefficient}\le C\tau)+R^u_j,
\]
$m=\beta_\sigma^{-1}(m')$, $\beta^u(m')=b(m',\sigma^s(m')+\eta^u(m'))$, where $R_j$ is a polynomial in $J_1,\dots,J_{j-1}$ with coefficients that are continuous functions of $(m,\sigma)$ built from derivatives of $g,b,P^{s,u}$ up to order $j$ (and for $j=1$ the equation for $\mu^u_1$ is solved for $\mu^u_1$, its own coefficient being $O(\tau)$). Hence
\[
\mathrm{Lip}_{J_j}\mu_j\le\max\{\|B^s\|\|D\beta^{-1}\|^j,\ \|(B^u)^{-1}\|\|D\beta^u\|^j\}+C\tau\le\max\{s_j,u_j\}+C\tau<1,
\]
\textbf{uniformly in $\sigma$}, and $\mu_j$ depends Lipschitz-continuously on $J_i$ ($i<j$) on the balls. $\Lambda$ is continuous in the product of sup metrics (uniform continuity of $J$ and of $\beta_\sigma^{-1}$ in $\sigma$, Section~5.3(c)). With a weighted fibre metric (Section~3.2) (FCT) applies. Start at $(\sigma_0,J^{(0)})=(0,0)$; as in Theorem 3.2(b), $\Lambda^n(0,0)=(\sigma_n,D\sigma_n,\dots,D^k\sigma_n)$ with $\sigma_{n+1}=\Gamma\sigma_n\in C^{k}$ (implicit function theorem for (5.2); $P^{s,u},B,g,b$ have the required regularity). So $\sigma_\varepsilon\in C^k$, $D^i\sigma_\varepsilon=J_i$.
\textbf{$C^{k,\alpha}$.} With $\sigma_\varepsilon$ fixed, $J_k$ is the fixed point of the fibre map $\mu_k(\sigma_\varepsilon,\cdot)$. By the one-step numbers $s_{k,\alpha},u_{k,\alpha}<1$ and $[J\circ\beta^{-1}]_\alpha\le[J]_\alpha\mathrm{Lip}(\beta^{-1})^\alpha$ etc., $[\mu_k(J)]_\alpha\le(\max\{s_{k,\alpha},u_{k,\alpha}\}+C\tau)[J]_\alpha+C(1+\|J\|_0)$ (for $k=1$ the Hölder seminorm of $D\beta^J$ is $\le C+K_*[J]_\alpha$, absorbed). Iterating from $J=0$ and passing to the uniform limit: $[J_k]_\alpha<\infty$. (The coefficients lie in $C^\alpha$ because $D_\eta g\in C^{r-1,\beta}$, $\sigma_\varepsilon$ is Lipschitz and $\alpha\le\beta$ if $k=r$.) $\beta_\varepsilon=b(\cdot,\sigma_\varepsilon)\in C^{k,\alpha}$ is a homeomorphism with $D\beta_\varepsilon$ invertible, hence a $C^{k,\alpha}$ diffeomorphism.

**(c)** At $(\varepsilon,\sigma,J)=(0,0,0)$ the lifted fibre maps have $R_j=0$ (all pure-$m$ derivatives of $g_0$ at $\eta=0$ vanish and $\sigma=0$ is invariant at $\varepsilon=0$), and $R_j$ is Lipschitz in $(\delta_\varepsilon^{(j)},\|\sigma\|_0,\|J_{<j}\|)$. Hence $\|R_j\|\le C(\delta^{(j)}_\varepsilon+\|\sigma\|_0+\sum_{i<j}\|J_i\|)$ and $\|J_j\|\le\|R_j\|/(1-\theta')$; induct, using $\|\sigma_\varepsilon\|_0\le C\delta_\varepsilon$. The Hölder estimate is the same with the Hölder seminorm of $R_k$. $\square$

\textbf{Theorem 5.4.} Let $X$ be a Banach space, $M\subset X$ a compact $C^{r+1,\beta}$ submanifold, $F_\varepsilon$ $C^{r+1,\beta}$ maps near $M$ with $F_0(M)=M$, $f=F_0|_M$ a diffeomorphism, and suppose the normal cocycle $B$ satisfies (H1), (B) with $k\ge1$, (H3), $k+\alpha\le r+\beta$. Choose the complement $E$ as in Lemma 3.4 with $\sup\|K\|\le\delta^*_1$. Then there is $\varepsilon_0>0$ such that for $|\varepsilon|<\varepsilon_0$: \textit{(flow reduction in Appendix G; regularity in Appendix A)}
\begin{itemize}
    \item[(a)] There is a unique $C^{k,\alpha}$ compact submanifold $M_\varepsilon=K_\varepsilon(M)$, $K_\varepsilon(m)=\Theta(m,\sigma_\varepsilon(m))$, in the tube with $F_\varepsilon(M_\varepsilon)=M_\varepsilon$ (exact equality) among graphs in the class $\Sigma_{\rho,L}$; $K_\varepsilon:M\to X$ is a $C^{k,\alpha}$ embedding and $f_\varepsilon:=\beta_\varepsilon$ is a $C^{k,\alpha}$ diffeomorphism with $F_\varepsilon\circ K_\varepsilon=K_\varepsilon\circ f_\varepsilon$.
    \item[(b)] If $\|F_\varepsilon-F_0\|_{C^{k,\alpha}}\le C|\varepsilon|$ then $\|K_\varepsilon-\iota_M\|_{C^{k,\alpha}}+\|f_\varepsilon-f\|_{C^{k,\alpha}}\le C'|\varepsilon|$.
    \item[(c)] $M_\varepsilon$ is normally hyperbolic: the normal cocycle $\tilde B_m=D_\eta g_\varepsilon(m,\sigma_\varepsilon(m))-D\sigma_\varepsilon(\beta_\varepsilon m)D_\eta b_\varepsilon(m,\sigma_\varepsilon(m))$ over $f_\varepsilon$ satisfies (H1) with constants $\lambda_\varepsilon\to\lambda$, and its bunching numbers satisfy $\bar{\mathfrak b}^\pm_{k,\alpha}(\varepsilon)\le\bar{\mathfrak b}^\pm_{k,\alpha}+o(1)$.
    \item[(d)] \textit{Flows.} If $\Phi^\varepsilon_t$ is a $C^{r+1,\beta}$ local flow with $\Phi^0_t(M)=M$ for $t\in[0,T]$ and $F_\varepsilon:=\Phi^\varepsilon_T$ satisfies the hypotheses, then $\Phi^\varepsilon_t(M_\varepsilon)=M_\varepsilon$ for all $t\in\mathbb R$.
\end{itemize}

\textit{Proof.} (a),(b): Theorem 5.3 with $\widehat F_\varepsilon=\Theta^{-1}F_\varepsilon\Theta$; $F_\varepsilon(K_\varepsilon m)=\Theta\widehat F_\varepsilon(m,\sigma_\varepsilon m)=\Theta(\beta_\varepsilon m,\sigma_\varepsilon(\beta_\varepsilon m))=K_\varepsilon(\beta_\varepsilon m)$ and $\beta_\varepsilon$ is bijective. (c) $T_{K_\varepsilon m}M_\varepsilon$ is spanned by $(v,D\sigma_\varepsilon v)$, so the normal action on the fibre $E$ is $\eta\mapsto D_\eta g\,\eta-D\sigma_\varepsilon(\beta m)D_\eta b\,\eta$. Its $N$-step blocks in the reference splitting differ from those of $B$ by $O(C_N\tau)$, so the one-step quantities $s_j,u_j$ change by $O(\tau)$ and (by Lemma 2.2(d), upper semicontinuity of the Fekete limits) so do the numbers $\bar{\mathfrak b}^\pm$; an invariant splitting exists by Theorem 3.2(a) applied to the Grassmannian fibre map as in Lemma 3.3. (d) For $t\in[0,T]$, $\Phi^\varepsilon_t(M_\varepsilon)$ is $F_\varepsilon$-invariant ($\Phi_t$ commutes with $\Phi_T$) and, being $C^1$-close to $\Phi^0_t(M)=M$ uniformly on $[0,T]$, is a graph over $M$ (Lemma 5.1 for its base projection) of a section in $\Sigma_{\rho,L}$. By uniqueness it equals $M_\varepsilon$. For other $t$ use $\Phi_{t+T}=\Phi_t\Phi_T$. $\square$

\textbf{Theorem 5.5.} Assume (H4). \textit{(Lipschitz/differentiability via parameter expansions, Appendix F; linearised cohomological equation inverted by Theorem 4.2/Appendix A)}
\begin{itemize}
    \item[(a)] \textit{Lipschitz dependence.} $\|\sigma_\varepsilon-\sigma_{\varepsilon'}\|_{C^j}\le C\|\widehat F_\varepsilon-\widehat F_{\varepsilon'}\|_{C^j}$ for $j\le k$.
    \item[(b)] \textit{Differentiability with loss.} If $\varepsilon\mapsto\widehat F_\varepsilon$ is $C^1$ into $C^{k+1}(E_{\rho_0})$ and $k\ge1$, then $\varepsilon\mapsto\sigma_\varepsilon\in C^{k-1}(E)$ is $C^1$, and $\xi=\partial_\varepsilon\sigma_\varepsilon$ is the unique solution of the linearised cohomological equation $\xi\circ\beta_\varepsilon-\tilde B\xi=-(\partial_\varepsilon\Psi)(\varepsilon,\sigma_\varepsilon)$, $\Psi(\varepsilon,\sigma)=\sigma\circ\beta_\sigma-g_\varepsilon(\cdot,\sigma)$. Inductively, if $\varepsilon\mapsto\widehat F_\varepsilon\in C^p(\,\cdot\,;C^{k+p}(E_{\rho_0}))$ and $p\le k$ then $\sigma_\varepsilon\in C^p(\,\cdot\,;C^{k-p})$ and the Taylor coefficients satisfy $\mathcal A_0\sigma^{(p)}=-R_p(\sigma^{(1)},\dots,\sigma^{(p-1)})$ with $R_p$ the universal Faà di Bruno expression.
\end{itemize}

\textit{Proof.} (a) The fixed points of the uniform contractions $\Gamma_\varepsilon$ and of the affine fibre maps $\mu_j$ depend Lipschitz on the parameter with constant $(1-\theta)^{-1}$ times the Lipschitz constant of the family (CT); the family distance is $\le C\|\widehat F_\varepsilon-\widehat F_{\varepsilon'}\|_{C^j}$ plus the distances of the lower jets; induct on $j$.
(b) The linearisation $D_\sigma\Psi(\varepsilon,\sigma_\varepsilon)\xi=\xi\circ\beta_\varepsilon-\tilde B\xi$ is conjugate to $I-D\Gamma$, with $\|D\Gamma\|\le\theta<1$ in $C^0$ and, by the jet estimates of Theorem 5.3(b), uniformly on $C^j$ ($j\le k-1$ suffices below); hence it is invertible there with bounded inverse. Let $D_h=(\sigma_{\varepsilon+h}-\sigma_\varepsilon)/h$. By (a), $\|D_h\|_{C^k}\le C$. Subtracting the invariance equations, $D_\sigma\Psi(\varepsilon,\sigma_\varepsilon)D_h=-\partial_\varepsilon\Psi+O(|h|)$ in $C^{k-1}$, where the $O(|h|)$ remainders are second-order differences in $(\sigma_{\varepsilon+h}-\sigma_\varepsilon,\beta_{\varepsilon+h}-\beta_\varepsilon)$, controlled in $C^{k-1}$ by $\|D_h\|_{C^k}|h|$ (one derivative is paid for the composition). Hence $D_h\to-(D_\sigma\Psi)^{-1}\partial_\varepsilon\Psi$ in $C^{k-1}$. Continuity of the derivative follows from (a) at order $k-1$. Higher derivatives by induction, losing one derivative each time. $\square$

\textbf{Theorem 5.6.} (a) \textit{Skew products.} Suppose $E=E^s$, $b_\varepsilon(m,\eta)=f(m)$ and $g_\varepsilon(m,\cdot)$ is a $C^{r,\beta}$ family of maps $E_{m}\to E_{fm}$ with $\mathfrak c_{0,0}<1$. Then Theorem 3.2 gives a unique invariant section, of class $C^{k,\alpha}$ whenever $\mathfrak c_{k,\alpha}<1$, \textbf{with no 1-bunching hypothesis}: in particular for $k=0$ and every $\alpha\in(0,1)$ with $\mathfrak c_{0,\alpha}<1$, even when $\mathfrak b_1\ge1$ and the section is merely Hölder. \textit{(Skew products and sharpness in Appendix K; see also Example 4.5(ii))}
(b) \textit{Sharpness.} In Example 4.5 take $F_\varepsilon(x,y)=(Ax,\lambda y+\varepsilon q(x))$ on $\mathbb T^2\times\mathbb R$, $q=\cos(2\pi\langle\nu,x\rangle)$. The invariant graph is $\sigma_\varepsilon=\varepsilon\,\mathcal Gq$, and by Example 4.5(ii) it lies in $C^s$ (non-integer $s$) iff $s<s_*=\log(1/\lambda)/\log\Lambda$, i.e. iff $\bar{\mathfrak b}_{k,\alpha}=\lambda\Lambda^{k+\alpha}<1$. So Theorems 5.3, 5.4, 5.5 are sharp: for every $s>s_*$ there are arbitrarily small smooth perturbations whose invariant section is not $C^s$.
(c) \textit{Where 1-bunching matters.} If $\mathfrak b_1^\pm<1$ fails the invariant set need not be a graph (it is a Hölder lamination in general); this is why Theorem 5.3 and Theorem 5.4 assume $\mathfrak b_1^\pm<1$, while (a) does not need it because the base map does not depend on the fibre variable.

% ============================================================================
% SECTION 6: STRONG STABLE FIBRES
% ============================================================================
\section{Strong stable fibres}

Let $F=F_\varepsilon$ be a $C^{r+1,\beta}$ map near the compact invariant manifold $M_\varepsilon$ (Theorem 5.4; $M_0=M$ for $\varepsilon=0$), $F|_{M_\varepsilon}$ a diffeomorphism, and for $p\in M_\varepsilon$ write $p_n=F^n(p)$, $A_p(n)=DF(p_n)\in L(X)$, with cocycle $\Phi_p(n,k)=A_p(n-1)\cdots A_p(k)$.

\textit{Appendix B: strong stable fibres / Perron constructions provide the $C^{k,\alpha}$ regularity theory below.}

\textbf{(G) (uniform strong gap).} $DF$ has a continuous invariant splitting $X=E^s_p\oplus E^{cu}_p$ along $M_\varepsilon$ ($TM_\varepsilon\subset E^{cu}$, bounded projections $P^s,P^{cu}$, $DF(p)E^{cu}_p=E^{cu}_{Fp}$, $DF(p)E^s_p\subset E^s_{Fp}$) and constants $C\ge1$, $0\le\beta_{cu}<\alpha_s$ with
\[
\|\Phi^s_p(n,k)\|\le Ce^{-\alpha_s(n-k)}\ (n\ge k),\qquad\|\Phi^{cu}_p(n,k)\|\le Ce^{\beta_{cu}(k-n)}\ (n\le k),
\]
where $\Phi^{cu}(n,k)=\Phi^{cu}(k,n)^{-1}$ for $n<k$.

\textit{Verification for $\varepsilon=0$.} By Lemma 3.4, $DF_0|_M$ has the invariant splitting $TM\oplus\mathcal N^s\oplus\mathcal N^u$; put $E^s=\mathcal N^s$, $E^{cu}=TM\oplus\mathcal N^u$. Then $\alpha_s=-\log\mathfrak b_0^+$ and $\beta_{cu}=\log\max\{1,L^+\}$ with $L^+=\lim(\ell_n^+)^{1/n}$ (the unstable normal part has backward rate $\le\lambda<1$); so (G) is the condition $\mathfrak b_0^+L^+<1$, which implies $\mathfrak b_1^+<1$. For small $\varepsilon$, (G) persists with $\alpha_s-o(1)$, $\beta_{cu}+o(1)$ (the $N$-step products are continuous in the data).

\textbf{Theorem 6.1 (strong stable fibres).} Assume (G), $\beta_{cu}<\gamma<\alpha_s$. There are $r_0,\delta>0$ (uniform in $p$) such that for every $p\in M_\varepsilon$:
\begin{itemize}
    \item[(a)] \textit{Existence and characterisation.} There is a map $w_p:\{a\in E^s_p:\|a\|\le\delta\}\to E^{cu}_p$, $\Xi_p(a)=p+a+w_p(a)$, such that a point $x=p+\xi_0$ with $\|\xi_0\|$ small satisfies $\sup_{n\ge0}e^{\gamma n}\|F^nx-p_n\|\le r_0$ \textit{iff} $x=\Xi_p(a)$ for some $a$, $\|a\|\le\delta$ (with $a=P^s_p\xi_0$); in that case $\|F^nx-p_n\|\le Ce^{-\gamma n}\|a\|$.
    \item[(b)] \textit{Tangency.} $w_p(0)=0$, $Dw_p(0)=0$; the leaf $W^{ss}_{loc}(p)=\Xi_p(B_\delta)$ is tangent to $E^s_p$.
    \item[(c)] \textit{Equivariance.} $F\big(W^{ss}_{loc}(p)\big)\subset W^{ss}_{loc}(F(p))$ after shrinking $\delta$ by the factor $e^{-\gamma}$.
    \item[(d)] \textit{Internal regularity.} $w_p\in C^{r+1,\beta}$ with norms bounded uniformly in $p$: \textbf{under (G) alone, with no bunching condition of the form $\alpha_s>\ell\beta_{cu}$.}
\end{itemize}

\textit{Proof.} Write $F^nx-p_n=\xi_n$; then $\xi_{n+1}=A_p(n)\xi_n+R_n(\xi_n)$, $R_n(\xi)=F(p_n+\xi)-F(p_n)-DF(p_n)\xi$, $\|R_n(\xi)\|\le C\|\xi\|^2$, $\|DR_n(\xi)\|\le C\|\xi\|$, uniformly in $p,n$. Let $\ell_\gamma=\{\xi=(\xi_n)_{n\ge0}:\|\xi\|_\gamma=\sup e^{\gamma n}\|\xi_n\|<\infty\}$ and
\[
\mathcal P(\xi;a)_n=\Phi^s_p(n,0)a+\sum_{k=0}^{n-1}\Phi^s_p(n,k{+}1)P^s_{p_{k+1}}R_k(\xi_k)-\sum_{k=n}^\infty\Phi^{cu}_p(n,k{+}1)P^{cu}_{p_{k+1}}R_k(\xi_k).
\]
\textit{Estimates.} If $\|\xi\|_\gamma\le r$: $e^{\gamma n}\|\Phi^s(n,0)a\|\le C\|a\|$; the stable sum is $\le C'r^2e^{\alpha_s}\sum_{k<n}e^{(\gamma-\alpha_s)n+(\alpha_s-2\gamma)k}\le C''r^2$: for $\alpha_s>2\gamma$ the sum is $\le Ce^{-\gamma n}$, for $\alpha_s=2\gamma$ it is $\le Cne^{-\gamma n}$, and for $\alpha_s<2\gamma$ it is $\le C$ (terms decrease in $k$ and $\gamma<\alpha_s$); the $cu$-sum is $\le e^{\gamma n}\sum_{j\ge0}Ce^{\beta_{cu}(j+1)}C'r^2e^{-2\gamma(n+j)}=C'r^2e^{\beta_{cu}}e^{-\gamma n}\sum_je^{-(2\gamma-\beta_{cu})j}<\infty$ because $2\gamma>\beta_{cu}$. The same estimates with $\|DR_n\|\le Cr e^{-\gamma n}$ show $\mathrm{Lip}(\mathcal P(\cdot;a))\le C r$. Choose $r$ with $Cr\le\frac12$ and $\delta=r/(2C)$: $\mathcal P(\cdot;a)$ is a $\frac12$-contraction of the ball $B_r\subset\ell_\gamma$ for $\|a\|\le\delta$. Let $\xi(a)$ be its fixed point; $\xi(0)=0$.
\begin{itemize}
    \item[(a)] If $\xi$ is the orbit deviation of some $x$ with $\|\xi\|_\gamma\le r$, project the recursion with the invariant splitting: $\xi^s_n=\Phi^s(n,0)\xi^s_0+\sum_{k<n}\Phi^s(n,k+1)P^sR_k$ and, for $N>n$, $\xi^{cu}_n=\Phi^{cu}(n,N)\xi^{cu}_N-\sum_{k=n}^{N-1}\Phi^{cu}(n,k+1)P^{cu}R_k$. Since $\|\Phi^{cu}(n,N)\xi^{cu}_N\|\le Ce^{\beta_{cu}(N-n)}re^{-\gamma N}\to0$ ($\beta_{cu}<\gamma$), $\xi=\mathcal P(\xi;P^s\xi_0)$, so $\xi=\xi(a)$ by uniqueness. Conversely $\xi(a)$ is an orbit deviation by construction. Define $w_p(a)=P^{cu}\xi_0(a)$; then $\xi_0(a)=a+w_p(a)$.
    \item[(b)] $\|w_p(a)\|\le C\sum_k\|R_k\|\le Cr^2$-type bound sharpened by $\|\xi(a)\|_\gamma\le C\|a\|$ to $\|w_p(a)\|\le C\|a\|^2$.
    \item[(c)] If $\xi=\xi(a)$ then $\xi'_n:=\xi_{n+1}$ solves the problem for $(Fp,a'=P^s\xi_1)$ and $\|\xi'\|_\gamma\le e^{-\gamma}\|\xi\|_\gamma$.
    \item[(d)] The Nemytskii map $\xi\mapsto(R_n(\xi_n))_n$ satisfies $\|D^i\mathcal N(\xi)(\eta_1,\dots,\eta_i)\|_{\ell_{\max\{2,i\}\gamma}}\le C\prod\|\eta_j\|_\gamma$ ($i\le r+1$) with $D^{r+1}$ uniformly continuous, and the two convolution operators map $\ell_{2\gamma}\to\ell_\gamma$ boundedly (above). Hence $\mathcal P\in C^{r+1}$ in $(\xi,a)$, and the fixed point of a $C^{r+1}$ uniform contraction is $C^{r+1}$ in the parameter. $\square$
\end{itemize}

\textbf{Remark 6.1' (non-reflexive Perron construction).} The half-line Perron construction above uses only the completeness of the weighted Banach sequence space $\ell_\gamma$ and the uniform contraction estimates; it does \textbf{not} require reflexivity of the ambient Banach space $X$, nor weak compactness of bounded sets. The same remark applies to the fibre contraction theorem of §3 and to the mixed graph transform of §5. This concerns internal leaf construction and internal regularity; transverse regularity of the family of leaves and the phase projection remain separate questions requiring additional bunching (Remark 6.2).

\textbf{Remark 6.2 (transverse regularity; conjectural threshold).} The regularity of $p\mapsto W^{ss}(p)$ (holonomy) is governed by \textit{forward} separation rates of base orbits, not by the normal jet cocycle: differentiating the Perron operator $j$ times in $p$ produces weights $\gamma-j\mu$ with $\mu$ the forward growth rate of $\|DF^n\|$ along $M_\varepsilon$, so the natural conjecture is $C^\ell$ transversally under $\alpha_s>\ell\mu+\beta_{cu}$ (compare Pugh–Shub–Wilkinson for Hölder exponents). A rigorous proof requires comparing leaves through different base points in a common trivialisation and a scale version of the uniform contraction principle; this is Problem 2 of Section 9.

% ============================================================================
% SECTION 7: PERIODIC JET RESONANCE
% ============================================================================
\section{Periodic jet resonance}

Let $p\in M$ be periodic, $f^N(p)=p$ with prime period $N$, $m_i=f^i(p)$. Let $J^k_x(E)=\bigoplus_{j\le k}\mathrm{Sym}^j(T^*_xM;E_x)$ (jets of sections at $x$ in a fixed chart and trivialisation). For $\zeta\in C^k(E)$ write $J^k\zeta(x)\in J^k_x(E)$. Differentiating $\mathcal A\zeta=q$ gives, at each orbit point,
\[
\mathsf T_{m_i}\,J^k\zeta(m_{i+1})-\mathsf B_{m_i}\,J^k\zeta(m_i)=J^kq(m_i),\tag{7.1}
\]
where $\mathsf T_x:J^k_{fx}(E)\to J^k_x(f^*E)$ is the jet-composition map $J\mapsto J\circ f$ at $x$ (invertible, because $f$ is a diffeomorphism) and $\mathsf B_x:J^k_x(E)\to J^k_x(f^*E)$ is jet-multiplication by $B$ (so $\mathsf T_x^{-1}\mathsf B_x:J^k_x(E)\to J^k_{fx}(E)$). The equations are block lower-triangular in the order $j$ with diagonal blocks $J\mapsto J\circ(Df)$ and $B_x\cdot$. Put $\mathsf R_i=\mathsf T_{m_i}^{-1}\mathsf B_{m_i}$ and $\mathfrak R_{\le k}(p)=\mathsf R_{N-1}\cdots\mathsf R_0$, a linear map of $J^k_p(E)$, block lower-triangular with diagonal blocks (the \textbf{principal jet return maps})
\[
\mathfrak R_j(p)J=\Phi(N,0;p)\,J\circ\big(Df^N(p)^{-1}\big)^{\otimes j}\qquad(J\in\mathrm{Sym}^j(T^*_pM;E_p)).
\]

\textbf{Proposition 7.1 (graded periodic obstruction).} Assume $\mathcal A$ is surjective as a map $C^{k,\alpha}(E)\to C^{k,\alpha}(f^*E)$ (or $C^k\to C^k$). Then $I-\mathfrak R_{\le k}(p)$ is surjective on $J^k_p(E)$ for every periodic $p$. Equivalently: if for some periodic $p$ the map $I-\mathfrak R_{\le k}(p)$ is not onto $J^k_p(E)$, then $\mathcal A$ is not onto on $C^k$ (nor $C^{k,\alpha}$).

\textit{Proof.} Given $Z\in J^k_p(E)$ choose $q\in C^{k,\alpha}$ (a polynomial in a chart times a bump) supported near $m_{N-1}$ with $\mathsf T_{m_{N-1}}^{-1}J^kq(m_{N-1})=Z$ and $J^kq(m_i)=0$ for $i\ne N-1$. If $\mathcal A\zeta=q$ with $\zeta\in C^k$, iterate $J^k\zeta(m_{i+1})=\mathsf R_iJ^k\zeta(m_i)+\mathsf T^{-1}_{m_i}J^kq(m_i)$ around the cycle: $Y=\mathfrak R_{\le k}Y+Z$ for $Y=J^k\zeta(p)$. So $Z\in\mathrm{ran}(I-\mathfrak R_{\le k})$. $\square$

\textbf{Example 7.2 (two-cycle).} $M=S^1$ with a diffeomorphism having an attracting 2-cycle $\{p_0,p_1\}$ with $f'(p_0)=f'(p_1)=\tfrac12$; $E=M\times\mathbb R$, $B\equiv\tfrac12$ ($E=E^s$, $\lambda=\tfrac12$). Order 0: $\mathfrak R_0=\Phi(2,0)=\tfrac14\ne1$, so $\mathcal A$ is invertible on $C^0$ (matrix $\begin{pmatrix}-1/2&1\\1&-1/2\end{pmatrix}$, determinant $-\tfrac34$). Order 1: $\mathfrak R_1=\Phi(2,0)(Df^2)^{-1}=\tfrac14\cdot4=1$, so $I-\mathfrak R_1=0$: $\mathcal A$ is \textbf{not onto $C^1$}. (Here $\mathfrak b_1^+\ge1$: the system sits exactly at the 1-bunching threshold.)

\textbf{Theorem 7.3 (compact return).} Assume additionally that $E^u_p$ is finite-dimensional and $\Phi^s(N,0;p)$ is compact on $E^s_p$ (as for eventually compact semiflows; then $B_x$ is non-invertible on $E^s$, which (H1) permits). Then $I-\mathfrak R_{\le k}(p)$ is Fredholm of index $0$; hence it fails to be onto iff $1$ is an eigenvalue of some diagonal block $\mathfrak R_j(p)$, $j\le k$, i.e. iff there are an eigenvalue $\mu$ of $\Phi(N,0;p)$ and eigenvalues $d_1,\dots,d_j$ (with repetition) of $Df^N(p)$ with $\mu=d_1\cdots d_j$. In particular such a resonance with $j\le k$ is an obstruction to surjectivity of $\mathcal A$ on $C^{k,\alpha}$. \textit{(See Appendix K for the fractional jet-spectrum resonance obstruction; this theorem complements the fractional result: Theorem 7.3 applies to integer $k$, Appendix K extends to fractional jet spectra).}

\textit{Proof.} Each diagonal block $\mathfrak R_j$ is compact on its $E^s$-part (compact $\Phi^s(N,0)$ composed with finite-dimensional tensor factors) and finite-dimensional on its $E^u$-part; the off-diagonal blocks are compact as well. So $\mathfrak R_{\le k}(p)$ is compact and $I-\mathfrak R_{\le k}(p)$ is Fredholm of index $0$ (Riesz–Schauder). For a block-triangular operator with diagonal blocks of the form $I-\text{(compact)}$, the index is the sum of the indices (all $0$), and invertibility holds iff each diagonal block is invertible; $I-\mathfrak R_j$ is non-invertible iff $1\in\sigma_p(\mathfrak R_j)$. On $\mathrm{Sym}^j(T^*_pM;E_p)$ the eigenvalues of $J\mapsto\Phi J(Df^{-N})^{\otimes j}$ are $\mu\,d_{i_1}^{-1}\cdots d_{i_j}^{-1}$. Apply Proposition 7.1. $\square$

% ============================================================================
% SECTION 8: EXTENSIONS
% ============================================================================
\section{Extensions}

The compact invertible theorem is the model case. The table summarises the replacements used in the extension settings; each extension requires its own theorem-level hypotheses.

\begin{table}[h]
\centering
\begin{tabular}{l|l|l|l}
\toprule
Setting & Linear inverse & Function space & Geometric / analytic hypothesis\\
\midrule
Noncompact / Banach-dim base (§8.1) & Green operator (uniform) & $C^{k,\alpha}_b$ & bounded geometry / uniform atlas (H3') \\
Nonautonomous / random (§8.2) & pathwise Green operator & tempered weighted & uniform / tempered dichotomy \\
Parabolic slow manifold (§8.3) & analytic semigroup Perron & $C_\gamma([0,\infty),X_\theta)$ & sectorial dichotomy (P1)–(P3) \\
Delay (§8.4) & history-space Green operator & $C([-\tau,0];\mathbb R^d)$ & eventual compactness (D0)–(D2) \\
Quasilinear (§8.5) & tame inverse / Nash–Moser & scale $X_s$ & tame + loss (Q1)–(Q2) \\
Manifolds with boundary (§8.6) & Green operator with boundary conditions & $C^k(M,E)$ & boundary-compatible hyperbolicity \\
\bottomrule
\end{tabular}
\caption{Extension settings and replacements}
\label{tab:extensions}
\end{table}


\subsection{Noncompact and Banach-dimensional bases}

\textbf{Theorem 8.1.} Let $M$ be a complete $C^{r+1,\beta}$ manifold, finite- or infinite-dimensional, with a \textit{uniform atlas} (charts of a fixed size with uniformly bounded $C^{r+1,\beta}$ transition maps), $f\in\mathrm{Diff}^{r+1,\beta}(M)$ with $\|f^{\pm1}\|_{C^{r+1,\beta}_b}<\infty$, $E\to M$ a $C^{r,\beta}$ Banach bundle with uniformly bounded trivialisations, $B\in C^{r,\beta}_b$ with (H1) holding with uniform constants, and (H3'): there is a $C^{r,\beta}_b$ reference splitting uniformly $\delta'$-close to $E^s\oplus E^u$ for every $\delta'>0$ (automatic for finite-dimensional $M$ of bounded geometry, by a uniformly locally finite partition of unity with uniformly bounded derivatives). Define all numbers with $\sup_{m\in M}$ and all function spaces as $C^{k,\alpha}_b$. Then Lemma 2.2, Lemma 3.1, Theorem 3.2, Lemma 3.3, Theorems 4.2–4.4, Proposition 5.2, and Theorems 5.3 and 5.5 hold verbatim (with constants depending only on the uniform bounds); the conclusion of Theorem 5.3(a) is a unique global invariant section $\sigma_\varepsilon\in C^{k,\alpha}_b$. \textit{(Global validity of the linear estimates via the noncompact uniform inverse of Appendix I; uniform constants via Appendix C (roughness) and Appendix F (parameter expansion).)}

\textit{Proof.} Compactness enters the proofs only through (i) Lemma 5.1 (proved above for complete $M$ with uniform charts, using no degree or covering argument); (ii) the existence of a finite atlas (replaced by the uniform atlas; every estimate in a chart is uniform); (iii) the partition-of-unity approximations of the projections in Lemma 3.3 and the complement in Lemma 3.4 (replaced by (H3') and the assumption that a complement with coupling $\le\delta_1^*$ exists); (iv) completeness of $C^0(V_\rho)$ and of the weighted fibre spaces, which holds for $C_b$-norms. No Arzelà–Ascoli argument is used. The fibre contraction theorem is stated for arbitrary complete metric spaces. $\square$

\subsection{Nonautonomous and random cocycles}

For nonautonomous and random extensions of normally hyperbolic persistence, see Latushkin–Montgomery-Smith–Randolph~\cite{lat1996} and Sell–You~\cite{sell2002}; the proof below adapts the time-independent theory to the bi-infinite setting.

\textbf{Theorem 8.2 (nonautonomous).} Let $(f_n,B_n,\widehat F_{\varepsilon,n})_{n\in\mathbb Z}$ be a bi-infinite sequence of data on the same $M,E$ (maps $f_n:M\to M$, fibre maps $B_n:E_m\to E_{f_nm}$, splittings $E^{s,u}_n$ with $B_nE^s_n\subset E^s_{n+1}$, $B_nE^u_n=E^u_{n+1}$) such that all one-step quantities
\[
\sup_m\|B^s_{n-1,\psi_{n-1}m}\|\kappa^+(m)^j\le\theta_0<1,\qquad\sup_m\|(B^u_{n,m})^{-1}\|\kappa^-(m)^j\le\theta_0<1\quad(j\le k,\ \text{with }\ell^\alpha\text{ at }j=k),
\]
and the Hölder norms of the data, are bounded uniformly in $n$, and $\sup_n\|\widehat F_{\varepsilon,n}-\widehat F_{0,n}\|_{C^{k,\alpha}}\to0$. Then for small $\varepsilon$ there is a unique bounded sequence of sections $\sigma_{\varepsilon,n}\in\Sigma_{\rho,L}$ with
\[
\sigma_{\varepsilon,n+1}\big(b_{\varepsilon,n}(m,\sigma_{\varepsilon,n}(m))\big)=g_{\varepsilon,n}\big(m,\sigma_{\varepsilon,n}(m)\big),
\]
with $\sigma_{\varepsilon,n}\in C^{k,\alpha}$ and the bounds of Theorem 5.3(c) uniform in $n$. \textit{(Bi-infinite sequence persistence; the one-step uniform estimates use Appendix A (linear inversion) and the parameter dependence uses Appendix F.)}

\textit{Proof.} On the complete space $\prod_n\Sigma_{\rho,L}$ (sup over $n$) let $(\Gamma\sigma)^s_{n+1}$ be defined from $\sigma_n$ by (5.1) with the data at time $n$, and $(\Gamma\sigma)^u_n$ from $\sigma_{n+1}$ by (5.2) with the data at time $n$. All estimates of Proposition 5.2 and of the jet lift are one-step and uniform in $n$, so the proofs of Theorem 5.3 (b), (c) go through verbatim (no iteration step is needed since the one-step numbers are already $<1$). $\square$

\textbf{Theorem 8.3 (random, uniformly hyperbolic).} Let $(\Omega,\mathcal F,\mathbb P,\theta)$ be an ergodic system and let $(f_\omega,B_\omega,\widehat F_{\varepsilon,\omega})$ be measurable random data (in the sense of Arnold~\cite{arn98}) satisfying the hypotheses of Theorem 8.2 with $n$ replaced by the time-shift $\omega\mapsto\theta\omega$ (uniformly in $\omega$). Then there is $\varepsilon_0>0$ and, for $|\varepsilon|<\varepsilon_0$, a \textbf{unique measurable family} $\sigma_\varepsilon(\omega)\in\Sigma_{\rho,L}$ with the \textit{cocycle} invariance
\[
\sigma_\varepsilon(\theta\omega)\Big(b_\varepsilon\big(\omega,m,\sigma_\varepsilon(\omega)(m)\big)\Big)=g_\varepsilon\big(\omega,m,\sigma_\varepsilon(\omega)(m)\big),
\]
i.e. $\widehat F_{\varepsilon}(\omega,\mathrm{Gr}\,\sigma_\varepsilon(\omega))=\mathrm{Gr}\,\sigma_\varepsilon(\theta\omega)$. \textit{(Random ergodic extension; see Appendix I for noncompact linear theory needed in measurable settings.)}

\textit{Proof.} Apply Theorem 8.2 to $n\mapsto(\ldots)_{\theta^n\omega}$ for each $\omega$; uniqueness identifies the sequence with the family $n\mapsto\sigma_\varepsilon(\theta^n\omega)$, and measurability follows because $\sigma_\varepsilon(\omega)=\lim_n\Gamma^n(0)(\omega)$ is a pointwise limit of measurable maps. $\square$

\subsection{Parabolic slow manifolds}

For the theory of invariant manifolds in infinite-dimensional settings via the parameterization method, see Haro et al.~\cite{har16}; the proof below follows the present fibre contraction approach.

\textbf{Hypotheses.} (P1) \textit{smooth dependence} For a fixed $T>0$ the time-$T$ map $F_\varepsilon=\Phi^\varepsilon_T$ of the semiflow on $X_\theta\times Y$ is $C^{r+1,\beta}_b$ near the critical manifold. (P2) \textit{critical manifold} $M_0=\{(h_0(v),v):v\in Y_0\}$ with $Ah_0(v)+f(h_0(v),v,0)=0$, $h_0\in C^{r+1,\beta}_b(Y_0;X_1)$, $Y_0\subset Y$ complete with a uniform atlas (e.g. $Y_0=Y$). (P3) \textit{uniform dichotomy} $L_v=A+D_uf(h_0(v),v,0)$ generates an analytic semigroup, and $\|e^{tL_v}|_{E^s_v}\|\le Ce^{-\eta t}$ ($t\ge0$), $\|e^{-tL_v}|_{E^u_v}\|\le Ce^{-\eta t}$ with $E^u_v$ finite-dimensional, uniformly in $v$, and the spectral projections $P^{s,u}_v$ (Riesz projections of $L_v$) are $C^{r,\beta}_b$ in $v$. (This follows from a uniform resolvent bound on the strip $\{|\mathrm{Re}\,z|\le\eta\}$ and uniform sectoriality; the hypothesis stated in v2, a resolvent bound on $\{\mathrm{Re}\,z\ge-\eta\}$, would exclude any unstable part.)

\textbf{Theorem 8.4.} Under (P1)–(P3) there are $\varepsilon_0>0$ and $h_\varepsilon\in C^{r,\beta}_b(Y_0;X_\theta)$, $|\varepsilon|<\varepsilon_0$, with $\|h_\varepsilon-h_0\|_{C^{r,\beta}_b}\le C|\varepsilon|$, such that $M_\varepsilon=\{(h_\varepsilon(v),v)\}$ is the unique invariant graph of $F_\varepsilon$ in a tube; it is invariant under the semiflow (for all $t\ge0$ for which solutions stay in the tube, and under the full flow of the reduced equation $v_t=\varepsilon g(h_\varepsilon(v),v,\varepsilon)$ on $M_\varepsilon$). \textit{(Parabolic critical-manifold section; see Appendix K (critical manifold jets / fractional spectra) and Appendix C (dichotomy roughness under the parabolic semigroup).)}

\textit{Proof.} Use the trivial bundle $E=Y_0\times X_\theta$ and $\Theta(v,\xi)=(h_0(v)+\xi,v)$, which is $C^{r+1,\beta}$. At $\varepsilon=0$ the $v$-component is constant, so $b_0(v,\xi)=v$: the base map is $f=\mathrm{id}$, the coupling $K\equiv0$, $\kappa^\pm_n\equiv1$, $\ell_n^\pm=1$. Also $g_0(v,0)=0$ (the critical manifold consists of equilibria) and $B_v=D_\xi g_0(v,0)=e^{TL_v}$, which satisfies (H1) with $a^\pm_n\le Ce^{-\eta nT}$ by (P3); $B_v$ is not invertible on $E^s$ but $B^u_v$ is an isomorphism of the finite-dimensional $E^u_v$. Therefore $\mathfrak b^\pm_s\le e^{-\eta T}<1$ and $\bar{\mathfrak b}^\pm_{k,\alpha}<1$ for \textbf{every} $k+\alpha\le r+\beta$ (bunching is automatic because $f=\mathrm{id}$). (H3) holds by (P3). Apply Theorem 8.1 (global version of Theorem 5.3) with $K=0$. Invariance of $M_\varepsilon$ under $\Phi^\varepsilon_t$, $t\in[0,T]$, follows as in Theorem 5.4(d): $\Phi^\varepsilon_t(M_\varepsilon)$ is an $F_\varepsilon$-invariant graph in the class (it is $C^1$-close to $\Phi_t^0(M_0)=M_0$, as $\Phi^0_t=\mathrm{id}$ on $M_0$), hence equals $M_\varepsilon$. $\square$

\subsection{Delay equations}

For the theory of invariant manifolds in functional differential equations, see Hale~\cite{hal77} and Sell–You~\cite{sell2002}; the proof below adapts the compact persistence theorem to delay systems.

Consider $\dot x(t)=G(x_t,\varepsilon)$ on $X=C([-\tau,0];\mathbb R^d)$ with $G\in C^{r+1,\beta}_b$. \textbf{Hypotheses.} (D0) the solution maps $T^\varepsilon_t$, $t>0$, are $C^{r+1,\beta}$ (standard for $G\in C^{r+1,\beta}$) and $T^\varepsilon_t$ is compact for $t\ge\tau$; (D1) $M_0\subset X$ is a compact finite-dimensional $C^{r+1,\beta}$ submanifold with $T^0_t(M_0)=M_0$ and $T^0_t|_{M_0}$ a $C^{r+1,\beta}$ flow ($t\in\mathbb R$); (D2) for some $T\ge\tau$ the normal cocycle $B_m$ induced by $DT^0_T(m)$ satisfies (H1) (necessarily with $E^u_m$ \textbf{finite-dimensional}, as $B$ is compact for $T\ge\tau$), (H3) and (B). See Hale~\cite{hal77} for the standard theory of functional differential equations.

\textbf{Theorem 8.5.} Under (D0)–(D2) there are $\varepsilon_0>0$ and, for $|\varepsilon|<\varepsilon_0$, a unique $C^{k,\alpha}$ compact submanifold $M_\varepsilon$ near $M_0$, invariant under $T^\varepsilon_t$ for all $t\ge0$, such that $T_t^\varepsilon|_{M_\varepsilon}$ extends to a $C^{k,\alpha}$ flow which is $C^1$-close to $T^0_t|_{M_0}$. \textit{(Delay equations; one-sided invariant graph via Appendix D.)}

\textit{Proof.} The time-$T$ map $F_\varepsilon=T^\varepsilon_T$ is $C^{r+1,\beta}$ and, restricted to $M_0$, a diffeomorphism; its normal cocycle satisfies (H1) with non-invertible stable part, which is allowed. Theorem 5.4 applies, and Theorem 5.4(d) (with $\Phi^0_t(M_0)=M_0$) gives invariance under the semiflow. $\square$

\subsection{Quasilinear problems (loss of derivatives)}

For numerical methods and applied perspectives on quasilinear and infinite-dimensional normally hyperbolic persistence, see Dellnitz–Junge~\cite{dell98}, Broer et al.~\cite{broe91}, and Kaper–Kaper~\cite{kape02}; the theoretical framework below addresses the loss-of-derivatives problem analytically.

\textbf{Theorem 8.6.} Let $\mathcal X=C^\infty(M;E)$ with the gradings $\|\cdot\|_s=\|\cdot\|_{C^s}$ and let $\Psi:(-\varepsilon_*,\varepsilon_*)\times U\to\mathcal X$, $U\subset\mathcal X$ a $C^0$-neighbourhood of $0$, $\Psi(\varepsilon,\sigma)=\sigma\circ\beta_\sigma-g_\varepsilon(\cdot,\sigma)$ (the invariance equation; for quasilinear $F_\varepsilon$ the map $\Psi$ may lose derivatives), satisfy:
\begin{itemize}
    \item[(Q1)] $\Psi$ is smooth and tame (in the sense of Hamilton, Part II.2) with $\Psi(0,0)=0$;
    \item[(Q2)] for $(\varepsilon,\sigma)$ near $(0,0)$ the linearisation $D_\sigma\Psi(\varepsilon,\sigma):\mathcal X\to\mathcal X$ is bijective with inverse $V\Psi(\varepsilon,\sigma,\cdot)$ smooth tame in $(\varepsilon,\sigma,q)$.
\end{itemize}
Then there are $\varepsilon_0>0$ and a smooth tame map $\varepsilon\mapsto\sigma_\varepsilon\in C^\infty(M;E)$, $|\varepsilon|<\varepsilon_0$, with $\Psi(\varepsilon,\sigma_\varepsilon)=0$.

\textit{Proof.} Apply the Nash–Moser theorem (Hamilton III.1.1.1) to $\mathcal P(\varepsilon,\sigma)=(\varepsilon,\Psi(\varepsilon,\sigma))$ at $(0,0)$: $D\mathcal P$ is bijective with smooth tame inverse $(\delta\varepsilon,q)\mapsto(\delta\varepsilon,V\Psi(\cdot,q-\partial_\varepsilon\Psi\,\delta\varepsilon))$. Hence $\mathcal P$ is locally a smooth tame diffeomorphism, and $\sigma_\varepsilon=\mathrm{pr}_2\mathcal P^{-1}(\varepsilon,0)$. $\square$

\textbf{Remark.} (Q2) is where the linear theory enters: when $\Psi$ loses no derivative, Theorems 4.2 and 4.3 give bounded inverses on every $C^{k,\alpha}$ with tame dependence on $\sigma$ through the constants $C_{\varepsilon}$ of Lemma 3.1, but Nash–Moser is then unnecessary because Theorems 5.3 and 5.5 apply directly. For genuinely quasilinear problems, (Q2) must be verified with a loss $d>0$; a general verification is not claimed. \textit{(Jet/cohomological estimates underlying (Q2) are developed in Appendix K (fractional jet spectra).)}

\subsection{Manifolds with boundary}

When \(M\) has non-empty boundary, ordinary invariance may fail: the graph transform requires the tube to be forward-invariant, and a neighbourhood of a boundary point need not be. Four notions of invariance are distinguished.

\begin{enumerate}
    \item \textbf{Local invariance} away from the boundary: \(F_\varepsilon(M_\varepsilon\cap\operatorname{int}M)=M_\varepsilon\cap\operatorname{int}M\).
    \item \textbf{Inflowing invariance}: the vector field points inward on \(\partial M\), so orbits enter the interior and the invariant manifold is locally invariant in the interior.
    \item \textbf{Overflowing invariance}: the vector field points outward, so orbits leave, and the invariant manifold is invariant in the exterior.
    \item \textbf{Cutoff invariance}: the vector field is cut off outside a prescribed region, so the invariant manifold is invariant for the modified system.
\end{enumerate}

Persistence in the boundary case requires \textbf{boundary-compatible hyperbolicity}: the normal hyperbolicity must hold uniformly up to the boundary, and the chosen invariance notion must be compatible with the boundary geometry. Under these hypotheses, the graph transform applies verbatim in the interior and up to the boundary, with the boundary conditions determining which of the four notions is realised. A full treatment requires the boundary-compatible tubular neighbourhood and the uniform estimates up to \(\partial M\); this is a genuine extension of the compact case, not a direct corollary. \textit{(One-sided invariant graph transforms relevant here are in Appendix D; the full boundary case is not claimed by the two-sided theorem 5.3 of Appendix J.)}

\subsection*{Supplementary Appendix Map}

The following table maps the main-text theorems to the supporting modules of the supplementary material.

\begin{table}[h]
\centering
\caption{Supplementary appendix map.}
\label{tab:supp-map}
\begin{tabular}{lll}
\textbf{Main result} & \textbf{Mechanism used} & \textbf{Appendix} \\
\hline
Thm 4.2 (isomorphism) & Banach-bundle cocycle inversion, Green series & A \\
Thm 4.3, 4.4 (spectral radii) & Gelfand, Fekete, Green & K \\
Thm 5.3 (existence, two-sided) & two-sided HPS graph transform & J \\
Thm 5.4 (regularity) & strong stable fibres, Perron contractions & B \\
Thm 5.5 (Lipschitz, loss) & parameter expansions, Faà di Bruno jets & F \\
Thm 5.6 (sharpness) & toral automorphism, jet spectra & K \\
Thm 6.1 (strong fibres) & Perron strong fibres & B \\
Thm 7.3 (compact return) & Fredholm/resonance & K \\
Thm 8.1 (global, noncompact) & noncompact uniform inverses & I \\
Thm 8.2, 8.3 (nonautonomous, random) & uniform one-step contractions & I \\
Thm 8.4 (parabolic) & critical manifold jets & K \\
Thm 8.5 (delay) & one-sided transform & D \\
Thm 8.6 (quasilinear, loss) & Nash–Moser / tame inverses & K \\
\end{tabular}
\end{table}


% ============================================================================
% SECTION 9: WHAT REMAINS OPEN
% ============================================================================
\section{What remains open}

\begin{enumerate}
    \item \textbf{Converse without (R), (RC).} Is $\rho^\pm_{k,\alpha}<1$ necessary for $\mathcal A$ to be an isomorphism of $C^{k,\alpha}$ in general? Equivalently, can the Green operator be bounded on $C^{k,\alpha}$ while the spectrum of $\mathcal S^\pm$ has components on both sides of the unit circle? Theorem 4.4(c) settles this under (R), (RC) — which hold, in particular, when the operator spectrum is a single annulus (one Lyapunov exponent) or, more generally, connected — and Example 4.5 is the rate-homogeneous toral case. The open case is a \textbf{disconnected} operator spectrum with $1$ in a gap; a counterexample requires a dominated splitting whose Hölder spectrum has a gap at the unit circle. A related question: the relation of $\rho^\pm_{k,\alpha}$ to the Mather~\cite{mather1968}/Sacker–Sell~\cite{sacker1978} spectrum of jet cocycles acting on Hölder spaces.

    \item \textbf{Transverse regularity of strong stable foliations.} Prove $C^{\ell,\alpha}$ regularity of the holonomy under the forward-rate condition of Remark 6.2; the threshold in terms of \textit{forward} growth.

    \item \textbf{Nonlinear sharpness beyond rate-homogeneous classes.} Theorem 5.6(b) proves sharpness in the toral class; a general statement requires a nonlinear analogue of Theorem 4.4(c).

    \item \textbf{Quasilinear tame inverses.} The maximal-regularity case (parabolic) verifies (Q2) with loss \(d=1\) (or \(d=2\)); the quasilinear hyperbolic case without smoothing remains open.

    \item \textbf{Tempered random theory.} A Pesin-type version of Theorem 8.3 with tempered tube radii.

    \item \textbf{Hölder regime without 1-bunching.} Theorem 5.7 proves $C^{0,\alpha}$ regularity for general $F_\varepsilon$ under the Hölder bunching $\mathfrak b_0(L^\pm)^\alpha<1$, with threshold $\alpha_*=\log(1/\mathfrak b_0)/\log L^+$. The sharp threshold (whether $\alpha_*$ is optimal) and the higher regularity $C^{k,\alpha}$ without 1-bunching remain open.

    \item \textbf{Manifolds with boundary.} Section 8.6 records the four invariance notions and the boundary-compatible hyperbolicity condition as a structural statement; a full proof requires the boundary-compatible tubular neighbourhood and the uniform estimates up to \(\partial M\). The sharp conditions on the boundary geometry and the choice of invariance notion remain to be worked out.

    \item \textbf{Self-contained two-sided Banach-bundle HPS theorem.} The two-sided invariant-section theorem of §5 and the full HPS structural input require, at present, the complete matching lemma, graph contraction, Lipschitz preservation, and jet/difference-quotient extension for general two-sided situations; a fully self-contained proof within the present Banach-bundle framework remains open (the imported structural input is quoted in §1.2).
\end{enumerate}

% ============================================================================
% GLOSSARY
% ============================================================================
\section*{Glossary}

\begin{tabular}{l|l}
Normal cocycle $B$ & Bundle map $B_m:E_m\to E_{f(m)}$ describing linearized normal dynamics \\
Cohomological operator $\mathcal A$ & $\mathcal A\zeta = \zeta\circ f - B\zeta$; solves linearized persistence \\
Transfer operator $\mathcal S^\pm$ & Stable/unstable propagation along orbits; $\mathcal S^+$ forward, $\mathcal S^-$ backward \\
Bunching number $\mathfrak b^\pm_s$ & Fekete limit of $\sup_m a_n^\pm \kappa_n^{\pm s}$; measures contraction vs. amplification \\
Hölder bunching $\bar{\mathfrak b}^\pm_{k,\alpha}$ & Mixed Hölder-bunching; $\mathfrak b^\pm_k(L^\pm)^\alpha$ for rate-homogeneous systems \\
Spectral radius $\rho^\pm_{k,\alpha}$ & Actual invertibility threshold on function spaces; $\underline{\mathfrak b}\le\rho\le\bar{\mathfrak b}$ \\
Green operator $\mathcal G$ & Explicit inverse $\mathcal Gq=\sum(\mathcal S^+)^iP^s\hat q-\sum(\mathcal S^-)^nP^u\hat q$ \\
Fibre contraction & Fixed-point argument on sections; regularity of invariant sections via jet lift \\
Mixed graph transform & Two-sided invariant section: push stable forward, pull unstable backward \\
1-bunching $\mathfrak b_1^\pm<1$ & Condition ensuring the invariant set is a graph (not just a lamination) \\
Principal jet obstruction & Periodic orbit resonance blocking surjectivity at order $k$ \\
\end{tabular}

% ============================================================================
% DECLARATIONS
% ============================================================================
\section*{Declarations}

\textbf{Author information.} A. Abaee (independent researcher, Tehran, Iran). ORCID: \href{https://orcid.org/0000-0002-0019-1842}{0000-0002-0019-1842}. Email: \href{mailto:amin_abaee@ut.ac.ir}{amin\_abaee@ut.ac.ir}.

\textbf{Funding.} None.

\textbf{Competing interests.} None.

\textbf{AI use.} Ernie (Baidu) and Deepseek AI assisted with computations, numerical checks and literature review. The author reviewed and edited the outputs and takes full responsibility for the work.

\textbf{Data availability.} No separate data set is associated with this manuscript; the mathematical content, numerical checks, and literature entries are reported in the manuscript body and appendices and are archived in the accompanying source files.

\textbf{Ethics.} This manuscript concerns mathematical analysis of dynamical systems and does not involve human or animal participants, their data, or clinical material; therefore ethical approval and the corresponding informed consent were not applicable.

% ============================================================================
% REFERENCES
% ============================================================================
\begin{thebibliography}{25}

\bibitem{tre2010} Treschev, D., \& Zubelevich, O. \textit{Introduction to the Theory of Invariant Manifolds}. Springer, 2010 (survey).

\bibitem{mather1968} Mather, J. N. \textit{Anosov diffeomorphisms and flat manifolds}. \textit{Topology}, 7(3), 319–348, 1968.

\bibitem{sacker1978} Sacker, R. J., \& Sell, G. R. \textit{A spectral theory for linear differential equations}. \textit{Journal of Differential Equations}, 30(2), 243–279, 1978.

\bibitem{bates1995} Bates, P. W., Lu, K., \& Zeng, C. \textit{Persistence of invariant manifolds for partial differential equations}. \textit{Journal of Differential Equations}, 117(2), 481–525, 1995.

\bibitem{hal77} Hale, J. K. \textit{Theory of Functional Differential Equations}. Springer, 1977.

\bibitem{lat1996} Latushkin, Y., Montgomery-Smith, S., \& Randolph, T. \textit{Evolutionary semigroups and dichotomy of linear skew-product flows on Banach spaces}. \textit{Journal of Differential Equations}, 125(1), 73–116, 1996.

\bibitem{sell2002} Sell, G. R., \& You, Y. \textit{Dynamics of Evolutionary Equations}. Springer, 2002.

\bibitem{dell98} Dellnitz, M., \& Junge, O. \textit{On the approximation of complicated dynamical behaviour. II}. \textit{SIAM Journal on Numerical Analysis}, 36(4), 1486–1498, 1999.

\bibitem{broe91} Broer, H. W., Dumortier, D., van Strien, S. J., \& Takens, F. \textit{Structures in Dynamics}. North-Holland, 1991.

\bibitem{kape02} Kaper, M. J. B., \& Kaper, T. J. \textit{Asymptotic analysis of two reduction methods for systems of chemical reactions}. \textit{Physical Review E}, 65(3), 031504, 2002.

\bibitem{har16} Haro, A., Canadell, M., Figueras, J.-L., Luque, A., \& Mondelo, J.-M. \textit{The Parameterization Method for Invariant Manifolds}. Springer, 2016.

\bibitem{arn98} Arnold, L. \textit{Random Dynamical Systems}. Springer, 1998.

\bibitem{hir77} Hirsch, M. W., Pugh, C. C., \& Shub, M. \textit{Invariant Manifolds}. Springer, 1977.

\bibitem{pugh1970} Pugh, C., \& Shub, M. \textit{Stable manifolds and hyperbolic sets}. In \textit{Global Analysis}, Proc. Sympos. Pure Math. 14, AMS, 1970, 133–163.

\bibitem{hal1982} Hamilton, R. S. \textit{The inverse function theorem of Nash and Moser}. \textit{Bulletin of the American Mathematical Society}, 7(4), 65–222, 1982.

\end{thebibliography}

% ============================================================================
% END OF DOCUMENT
% ============================================================================
\end{document}
